% Options for packages loaded elsewhere
% Options for packages loaded elsewhere
\PassOptionsToPackage{unicode}{hyperref}
\PassOptionsToPackage{hyphens}{url}
%
\documentclass[
  letterpaper,
  11pt,
  twoside,
  openany]{scrbook}
\usepackage{xcolor}
\usepackage[inner=1.2in,outer=1.0in,top=1.0in,bottom=1.0in]{geometry}
\usepackage{amsmath,amssymb}
\setcounter{secnumdepth}{5}
\usepackage{iftex}
\ifPDFTeX
  \usepackage[T1]{fontenc}
  \usepackage[utf8]{inputenc}
  \usepackage{textcomp} % provide euro and other symbols
\else % if luatex or xetex
  \usepackage{unicode-math} % this also loads fontspec
  \defaultfontfeatures{Scale=MatchLowercase}
  \defaultfontfeatures[\rmfamily]{Ligatures=TeX,Scale=1}
\fi
\usepackage{lmodern}
\ifPDFTeX\else
  % xetex/luatex font selection
\fi
% Use upquote if available, for straight quotes in verbatim environments
\IfFileExists{upquote.sty}{\usepackage{upquote}}{}
\IfFileExists{microtype.sty}{% use microtype if available
  \usepackage[]{microtype}
  \UseMicrotypeSet[protrusion]{basicmath} % disable protrusion for tt fonts
}{}
\makeatletter
\@ifundefined{KOMAClassName}{% if non-KOMA class
  \IfFileExists{parskip.sty}{%
    \usepackage{parskip}
  }{% else
    \setlength{\parindent}{0pt}
    \setlength{\parskip}{6pt plus 2pt minus 1pt}}
}{% if KOMA class
  \KOMAoptions{parskip=half}}
\makeatother
% Make \paragraph and \subparagraph free-standing
\makeatletter
\ifx\paragraph\undefined\else
  \let\oldparagraph\paragraph
  \renewcommand{\paragraph}{
    \@ifstar
      \xxxParagraphStar
      \xxxParagraphNoStar
  }
  \newcommand{\xxxParagraphStar}[1]{\oldparagraph*{#1}\mbox{}}
  \newcommand{\xxxParagraphNoStar}[1]{\oldparagraph{#1}\mbox{}}
\fi
\ifx\subparagraph\undefined\else
  \let\oldsubparagraph\subparagraph
  \renewcommand{\subparagraph}{
    \@ifstar
      \xxxSubParagraphStar
      \xxxSubParagraphNoStar
  }
  \newcommand{\xxxSubParagraphStar}[1]{\oldsubparagraph*{#1}\mbox{}}
  \newcommand{\xxxSubParagraphNoStar}[1]{\oldsubparagraph{#1}\mbox{}}
\fi
\makeatother


\usepackage{longtable,booktabs,array}
\newcounter{none} % for unnumbered tables
\usepackage{calc} % for calculating minipage widths
% Correct order of tables after \paragraph or \subparagraph
\usepackage{etoolbox}
\makeatletter
\patchcmd\longtable{\par}{\if@noskipsec\mbox{}\fi\par}{}{}
\makeatother
% Allow footnotes in longtable head/foot
\IfFileExists{footnotehyper.sty}{\usepackage{footnotehyper}}{\usepackage{footnote}}
\makesavenoteenv{longtable}
\usepackage{graphicx}
\makeatletter
\newsavebox\pandoc@box
\newcommand*\pandocbounded[1]{% scales image to fit in text height/width
  \sbox\pandoc@box{#1}%
  \Gscale@div\@tempa{\textheight}{\dimexpr\ht\pandoc@box+\dp\pandoc@box\relax}%
  \Gscale@div\@tempb{\linewidth}{\wd\pandoc@box}%
  \ifdim\@tempb\p@<\@tempa\p@\let\@tempa\@tempb\fi% select the smaller of both
  \ifdim\@tempa\p@<\p@\scalebox{\@tempa}{\usebox\pandoc@box}%
  \else\usebox{\pandoc@box}%
  \fi%
}
% Set default figure placement to htbp
\def\fps@figure{htbp}
\makeatother





\setlength{\emergencystretch}{3em} % prevent overfull lines

\providecommand{\tightlist}{%
  \setlength{\itemsep}{0pt}\setlength{\parskip}{0pt}}



 


% Quarto PDF preamble for Pictures of Inference
% Quarto manages: documentclass, geometry, fonts (default), graphicx,
% hyperref, microtype, captions, listings, syntax highlighting.
% This file adds only what the book needs on top.

% --- Math fonts and packages ---
\usepackage{mathpazo}
\usepackage{mathtools}
\usepackage{amssymb}
\usepackage{bm}

% --- Callout boxes ---
\usepackage[breakable,skins]{tcolorbox}

\newtcolorbox{keyinsight}{
  enhanced, breakable,
  colframe=orange!70!black, colback=orange!5,
  fonttitle=\sffamily\bfseries\small, title=Key Insight,
  boxrule=0.8pt, left=10pt, right=10pt, top=8pt, bottom=8pt,
  before skip=12pt, after skip=12pt,
}

\newtcolorbox{tryitnow}{
  enhanced, breakable,
  colframe=blue!50!black, colback=blue!3,
  fonttitle=\sffamily\bfseries\small, title=Try It Now,
  boxrule=0.8pt, left=10pt, right=10pt, top=8pt, bottom=8pt,
  before skip=12pt, after skip=12pt,
}

\newtcolorbox{inthemath}{
  enhanced, breakable,
  colframe=yellow!50!black, colback=yellow!4,
  fonttitle=\sffamily\bfseries\small, title=Inside the Math,
  boxrule=0.8pt, left=10pt, right=10pt, top=8pt, bottom=8pt,
  before skip=12pt, after skip=12pt,
}

\newtcolorbox{codelab}{
  enhanced, breakable,
  colframe=green!50!black, colback=green!3,
  fonttitle=\sffamily\bfseries\small, title=Code Lab,
  boxrule=0.8pt, left=10pt, right=10pt, top=8pt, bottom=8pt,
  before skip=12pt, after skip=12pt,
}

\newtcolorbox{expedition}{
  enhanced, breakable,
  colframe=teal!60!black, colback=teal!4,
  fonttitle=\sffamily\bfseries\small, title=Expedition,
  boxrule=0.8pt, left=10pt, right=10pt, top=8pt, bottom=8pt,
  before skip=12pt, after skip=12pt,
}

\newtcolorbox{frontiernote}{
  enhanced, breakable,
  colframe=purple!50!black, colback=purple!3,
  fonttitle=\sffamily\bfseries\small, title=Frontier Note,
  boxrule=0.8pt, left=10pt, right=10pt, top=8pt, bottom=8pt,
  before skip=12pt, after skip=12pt,
}

% --- Custom math commands ---
\newcommand{\E}{\operatorname{\mathbb{E}}}
\newcommand{\Var}{\operatorname{Var}}
\newcommand{\Cov}{\operatorname{Cov}}
\newcommand{\Corr}{\operatorname{Corr}}
\newcommand{\bx}{\bm{x}}
\newcommand{\by}{\bm{y}}
\newcommand{\bbeta}{\bm{\beta}}
\newcommand{\bvarepsilon}{\bm{\varepsilon}}
\newcommand{\bX}{\bm{X}}
\newcommand{\bY}{\bm{Y}}
\newcommand{\bI}{\bm{I}}

% Read-aloud aid for first-appearance notation
\newcommand{\readas}[2]{#1\;\textit{(read: #2)}}
\makeatletter
\@ifpackageloaded{tcolorbox}{}{\usepackage[skins,breakable]{tcolorbox}}
\@ifpackageloaded{fontawesome5}{}{\usepackage{fontawesome5}}
\definecolor{quarto-callout-color}{HTML}{909090}
\definecolor{quarto-callout-note-color}{HTML}{0758E5}
\definecolor{quarto-callout-important-color}{HTML}{CC1914}
\definecolor{quarto-callout-warning-color}{HTML}{EB9113}
\definecolor{quarto-callout-tip-color}{HTML}{00A047}
\definecolor{quarto-callout-caution-color}{HTML}{FC5300}
\definecolor{quarto-callout-color-frame}{HTML}{acacac}
\definecolor{quarto-callout-note-color-frame}{HTML}{4582ec}
\definecolor{quarto-callout-important-color-frame}{HTML}{d9534f}
\definecolor{quarto-callout-warning-color-frame}{HTML}{f0ad4e}
\definecolor{quarto-callout-tip-color-frame}{HTML}{02b875}
\definecolor{quarto-callout-caution-color-frame}{HTML}{fd7e14}
\makeatother
\makeatletter
\@ifpackageloaded{bookmark}{}{\usepackage{bookmark}}
\makeatother
\makeatletter
\@ifpackageloaded{caption}{}{\usepackage{caption}}
\AtBeginDocument{%
\ifdefined\contentsname
  \renewcommand*\contentsname{Table of contents}
\else
  \newcommand\contentsname{Table of contents}
\fi
\ifdefined\listfigurename
  \renewcommand*\listfigurename{List of Figures}
\else
  \newcommand\listfigurename{List of Figures}
\fi
\ifdefined\listtablename
  \renewcommand*\listtablename{List of Tables}
\else
  \newcommand\listtablename{List of Tables}
\fi
\ifdefined\figurename
  \renewcommand*\figurename{Figure}
\else
  \newcommand\figurename{Figure}
\fi
\ifdefined\tablename
  \renewcommand*\tablename{Table}
\else
  \newcommand\tablename{Table}
\fi
}
\@ifpackageloaded{float}{}{\usepackage{float}}
\floatstyle{ruled}
\@ifundefined{c@chapter}{\newfloat{codelisting}{h}{lop}}{\newfloat{codelisting}{h}{lop}[chapter]}
\floatname{codelisting}{Listing}
\newcommand*\listoflistings{\listof{codelisting}{List of Listings}}
\makeatother
\makeatletter
\makeatother
\makeatletter
\@ifpackageloaded{caption}{}{\usepackage{caption}}
\@ifpackageloaded{subcaption}{}{\usepackage{subcaption}}
\makeatother
\makeatletter
\@ifpackageloaded{sidenotes}{}{\usepackage{sidenotes}}
\@ifpackageloaded{marginnote}{}{\usepackage{marginnote}}
\makeatother
\usepackage{bookmark}
\IfFileExists{xurl.sty}{\usepackage{xurl}}{} % add URL line breaks if available
\urlstyle{same}
\hypersetup{
  pdftitle={Pictures of Inference},
  pdfauthor={Dr.~Ian Helfrich; Dr.~Elizaveta Gonchar},
  hidelinks,
  pdfcreator={LaTeX via pandoc}}


\title{Pictures of Inference}
\usepackage{etoolbox}
\makeatletter
\providecommand{\subtitle}[1]{% add subtitle to \maketitle
  \apptocmd{\@title}{\par {\large #1 \par}}{}{}
}
\makeatother
\subtitle{A Visual Course in Statistics, Calculus, and Econometrics}
\author{Dr.~Ian Helfrich \and Dr.~Elizaveta Gonchar}
\date{May 2026}
\begin{document}
\frontmatter
\maketitle

\renewcommand*\contentsname{Table of contents}
{
\setcounter{tocdepth}{2}
\tableofcontents
}

\mainmatter
\bookmarksetup{startatroot}

\chapter*{Preface}\label{preface}
\addcontentsline{toc}{chapter}{Preface}

\markboth{Preface}{Preface}

\emph{Working draft. Architecture committed; chapters arrive as they
finish.}

This is \emph{Pictures of Inference}. A visual course in statistics,
calculus, and econometrics, by Dr.~Ian Helfrich and Dr.~Elizaveta
Gonchar.

The book is built to take a reader from ``what is a correlation'' to the
working knowledge needed to read, evaluate, and contribute to research
at the frontier of empirical economics. Four parts: How to See, How to
Speak, How to Think, How to Discover. Visualization is primary. Math is
taught as a language. Multiple data types are treated as first-class.

For collaborators: see \texttt{docs/process.md} for how we build,
\texttt{docs/voice/charter.md} for what the narrator does and does not
do, and \texttt{docs/table-of-contents.md} for the chapter map.

For readers: start with chapter 1
(\href{prose/part1/01-numbers-on-the-page.qmd}{Numbers on the page}) and
read forward. Each chapter has a narrative layer that goes straight
through and requires no math beyond what came before. Underneath that,
technical sections develop the derivations and proofs for readers who
want them. At the end, a short frontier note points at where current
research is arguing. Read the layers that suit you, and come back for
the others later.

\part{How to See}

\chapter{Numbers on the page}\label{numbers-on-the-page}

\section{What you just did}\label{what-you-just-did}

Here is a small piece of a real-shaped dataset: bilateral trade flows
for six countries observed over a decade, one row per pair-of-countries
per year. Look at it for ten seconds before reading on.

{\def\LTcaptype{none} % do not increment counter
\begin{longtable}[]{@{}
  >{\raggedright\arraybackslash}p{(\linewidth - 16\tabcolsep) * \real{0.0426}}
  >{\raggedright\arraybackslash}p{(\linewidth - 16\tabcolsep) * \real{0.0709}}
  >{\raggedright\arraybackslash}p{(\linewidth - 16\tabcolsep) * \real{0.0709}}
  >{\raggedleft\arraybackslash}p{(\linewidth - 16\tabcolsep) * \real{0.1702}}
  >{\raggedleft\arraybackslash}p{(\linewidth - 16\tabcolsep) * \real{0.1560}}
  >{\raggedleft\arraybackslash}p{(\linewidth - 16\tabcolsep) * \real{0.1560}}
  >{\raggedleft\arraybackslash}p{(\linewidth - 16\tabcolsep) * \real{0.1064}}
  >{\centering\arraybackslash}p{(\linewidth - 16\tabcolsep) * \real{0.1206}}
  >{\centering\arraybackslash}p{(\linewidth - 16\tabcolsep) * \real{0.1064}}@{}}
\toprule\noalign{}
\begin{minipage}[b]{\linewidth}\raggedright
year
\end{minipage} & \begin{minipage}[b]{\linewidth}\raggedright
exporter
\end{minipage} & \begin{minipage}[b]{\linewidth}\raggedright
importer
\end{minipage} & \begin{minipage}[b]{\linewidth}\raggedleft
exports (USD millions)
\end{minipage} & \begin{minipage}[b]{\linewidth}\raggedleft
exporter GDP (USD T)
\end{minipage} & \begin{minipage}[b]{\linewidth}\raggedleft
importer GDP (USD T)
\end{minipage} & \begin{minipage}[b]{\linewidth}\raggedleft
distance (km)
\end{minipage} & \begin{minipage}[b]{\linewidth}\centering
common language
\end{minipage} & \begin{minipage}[b]{\linewidth}\centering
common border
\end{minipage} \\
\midrule\noalign{}
\endhead
\bottomrule\noalign{}
\endlastfoot
2023 & USA & CHN & 85,058 & 27.36 & 17.79 & 11,587 & 0 & 0 \\
2023 & USA & DEU & 28,039 & 27.36 & 4.46 & 7,850 & 0 & 0 \\
2023 & USA & JPN & 27,112 & 27.36 & 4.21 & 9,997 & 0 & 0 \\
2023 & USA & GBR & 34,137 & 27.36 & 3.34 & 6,979 & 1 & 0 \\
2023 & USA & KOR & 9,602 & 27.36 & 1.71 & 10,534 & 0 & 0 \\
2023 & CHN & USA & 145,847 & 17.79 & 27.36 & 11,587 & 0 & 0 \\
\end{longtable}
}

Ten seconds. What did your eye do?

It found the rows and the columns. It noticed each row contained the
same kinds of things in the same order. It registered that some columns
held names (USA, CHN, DEU) and some held numbers, and that the numbers
came in different magnitudes. It parsed ``85,058'' as eighty-five
thousand of something, possibly more if you bothered to figure out the
units, possibly the size of an entire country's exports to another
country in a year. You probably started asking questions. Why does China
export almost twice as much to the United States as the United States
exports back? Why does the US export more to the much smaller economies
of Germany and Japan than to Korea? What does ``common language'' mean
as a number rather than as a word, and why is it 1 for the US-UK row and
0 for everyone else?

Nobody taught you to do any of that. You just did it.

This book is about the difference between doing it instinctively and
doing it deliberately. Most people who say they are bad at math, or at
any of the disciplines that wear quantitative clothes, are not bad at
the seeing. They are bad at the next step, the one where what your eye
noticed gets written down precisely enough to argue with. Here is the
claim the rest of this book will try to demonstrate: the gap between
those two skills is much smaller than it looks, and bridging it does not
require you to become a different kind of person. You already do most of
the work. We are going to slow down and notice how.

Chapter 1 is the easiest one. It establishes a single fact about almost
every dataset you will ever meet, that data is rectangular and the
rectangle is not a notational accident, and then catalogs the kinds of
things that can live in a column. By the end of this chapter you will
have a vocabulary for talking about any dataset, the way you already
have one for talking about a paragraph of text. The vocabulary lets you
decompose a dataset into parts, ask what each part is doing, and figure
out which questions are worth asking before any analysis happens. The
trade slice on this page is going to be a recurring character. By
chapter five you will know it as well as you know your own kitchen.

A note on how to read this book before we start. You can read at three
depths. The narrative goes straight through and requires no math beyond
what came before. If you want the math, the technical sections develop
it properly without interrupting the narrative. At the end of every
chapter, a short frontier note points at what current research is
fighting about. Pick the depths that suit you, and come back for the
others later.

One more thing. When you encounter an equation, read it aloud. Out loud.
Even if you do not yet know what every symbol means. Equations are
sentences with a subject, a verb, and modifiers, and pretending they are
not is the single biggest mistake students make in their first
statistics course. We are going to fix that one immediately.

\section{The shape}\label{the-shape}

Every dataset you will ever look at is a rectangle. Some rectangles are
tiny, two rows by three columns. Some are vast, millions of rows by
hundreds of columns. The rule is the same at any size. Each row is one
thing being observed. Each column is one thing being measured about it.

In our trade slice, every row is one country's exports to one other
country in one specific year. The first row, USA to CHN in 2023, is one
observation. Just one. The next row is a different observation, USA to
DEU in 2023, with its own values for everything. The columns describe
the attributes you can attach to a row: the year, the exporter, the
importer, the dollar value of exports, the GDP of each side, the
great-circle distance between them, whether they share a language,
whether they share a border.

This sounds obvious. It is not. The first question to ask of any
dataset, before any analysis, is: \emph{what is one row?} Until you
answer that, you cannot do anything correctly with the data, because the
same numbers can mean wildly different things depending on what the row
represents.

Take our slice. One row is one country-pair-year. Three hundred rows
total, six countries with five trading partners each over ten years.
Suppose we instead aggregated to one-country-year. We would have sixty
rows (six countries times ten years), and ``exports'' would mean each
country's total exports to the world rather than to a specific partner.
Aggregate further to one-country and we would have six rows. Summarize
to one-year and we would have ten. Same underlying facts, four different
shapes, four different statistical worlds.

The unit of analysis is not a notational detail. It is a substantive
choice that decides everything downstream. A regression run on
country-year data answers different questions than a regression run on
dyad-year data, and the same coefficient labeled ``the effect of
distance on trade'' means almost incomparable things in the two
settings. When you read a paper, the first thing to figure out is what
one row of their dataset is. When you write your own analysis, the first
thing to decide is what you want one row to be.

A column, on the other hand, is one variable measured across all the
rows. The ``distance\_km'' column has the same definition in every row,
namely how far apart the two countries are. In our slice that number is
the same every time the same pair appears, which is itself a clue about
the data. Distance does not depend on the year. It is redundant across
the time dimension. (We will come back to that. It matters for how
regression standard errors should behave when observations within a
country-pair are correlated, which they are.)

\begin{figure}

\centering{

\pandocbounded{\includegraphics[keepaspectratio]{prose/part1/../../pictures/figures/output/fig_data_rectangle.pdf}}

}

\caption{\label{fig-rect}The fundamental shape of a dataset. The
horizontal stripe marks one row, which is one observation, in this case
one country-pair-year. The vertical stripe marks one column, which is
one variable measured across every observation. Whatever the size of the
rectangle, the meaning of rows and columns does not change.}

\end{figure}%

Burn this image into your visual cortex now. You will need it for the
rest of the book. Whatever you are looking at, however many variables
and however many observations, the rectangle is the underlying shape.
The skill the book teaches is what to do once you see it.

\section{What lives in a column}\label{what-lives-in-a-column}

Columns come in flavors. The flavor matters because it determines what
arithmetic the column will support.

The simplest distinction is between columns of \textbf{numbers} and
columns of \textbf{labels}. Dollar value of exports: number. Country
code USA: label. You can take the average of a number column and have it
mean something. You cannot take the average of a column of country
codes; the ``average country'' between USA and DEU is not a sensible
thing. The exports-weighted average distance between trade partners, by
contrast, is a real quantity that some people compute on purpose.

Statisticians divide variables into four flavors, and the difference
between them comes down to what operations the values support.

\textbf{Categorical, unordered (nominal).} The values are names with no
built-in ranking: country codes, sex, eye color, whether someone owns a
passport. The only sensible operations are counting how often each value
appears and asking whether two rows share the same value. USA and DEU
are different. There is no answer to ``is USA bigger than DEU.''

\textbf{Categorical, ordered (ordinal).} Order matters, but the gaps
between adjacent values do not. Education levels (high school, college,
graduate). Survey responses from strongly-disagree to strongly-agree.
Movie ratings out of five. The difference between ``agree'' and
``strongly agree'' is not necessarily the same size as the difference
between ``neutral'' and ``agree.'' You can sort the values. You can talk
about medians. Treating the numbers as if they were dollars and
computing an average is a substantive choice that needs defending in
print.

\textbf{Quantitative, discrete.} Counts of things. The number of
children in a household, the number of trade partners a country has in a
given year. The values are real numbers and arithmetic works, but the
values can only be whole numbers. You cannot have 2.4 children. The
Poisson distribution, which we will meet in chapter three, is the
natural home of this kind of variable.

\textbf{Quantitative, continuous.} Dollars, distances, time intervals,
temperatures. Arithmetic works fully. The exports column in our slice is
continuous, and so is GDP, and so is distance.

A common trap, especially when you are reading other people's data:
things that look like numbers and aren't. ZIP codes. Phone numbers.
Survey response IDs. The numeric country code 840 used for the United
States in some bilateral databases. These are stored as integers, and a
careless analyst will compute the mean ZIP code of a sample, get a
result that is mathematically valid and substantively meaningless, and
not notice. The number is a label that happens to look like a number.
The statistical operation has to match what the column actually
represents, not what its storage type happens to be.

Now the opposite trap: things that don't look like numbers and are.
Yes/no answers, encoded as 0 and 1. Dummy variables for whether a row
meets some condition. The ``common\_language'' column in our slice is
one of these. It is a yes/no, encoded as 0 or 1, and you can absolutely
compute its average. The mean of a 0/1 column is the proportion of rows
where the value is 1. We use this constantly.

\begin{figure}

\centering{

\pandocbounded{\includegraphics[keepaspectratio]{prose/part1/../../pictures/figures/output/fig_variable_types.pdf}}

}

\caption{\label{fig-types}A small taxonomy of variable types.
Quantitative variables come in continuous and discrete varieties;
categorical variables come in ordered (ordinal) and unordered (nominal)
varieties. Examples on the right are the kinds of columns you actually
meet in working datasets.}

\end{figure}%

You will spend more of your statistical life than you would guess just
figuring out what kind of variable each column is. It is the foundation
for everything that follows, and it is almost never written down
explicitly in the dataset itself. Doing it carefully is one of the easy
wins of getting good at this.

\section{First questions}\label{first-questions}

Every time you sit down with a new dataset, before any analysis, run a
short list of questions through your head. The list takes ten minutes
for a small dataset and an afternoon for a serious one. It is the
difference between an analysis you can defend and an analysis that
quietly breaks underneath you.

Eight questions. Memorize the list.

\textbf{1. What is one row?} We did this already. Until you can name the
unit of analysis without thinking, you cannot do anything else
correctly. In our trade slice, one row is one country-pair-year.

\textbf{2. How many rows, how many columns?} For our slice, 300 rows and
9 columns. This sounds like trivia. It is not. The shape sets the cost
of mistakes. With 300 rows you can eyeball almost anything; with 300
million you cannot. You will adopt different tools, do different sanity
checks, and worry about different things as the row count grows.

\textbf{3. What is the type of each column?} Number or label. If number,
continuous or discrete. If label, ordered or unordered. Walk every
column. Be specific about borderline cases (the
\texttt{common\_language} column is technically categorical but encoded
as 0/1 so you can average it, and you will).

\textbf{4. What are the units?} ``Exports'' is meaningless without ``USD
millions, current prices.'' ``Distance'' needs ``kilometers.'' ``GDP''
needs ``trillions of USD'' and ideally ``nominal vs.~real, base year.''
A staggering fraction of bad analyses come from someone forgetting that
the units in two columns don't match. The codebook in
\texttt{data/trade/codebook.md} is the source of truth for our slice.
Read it.

\textbf{5. What is the time period? What is the geographic scope?} Our
slice covers 2014 to 2023, six countries chosen for size and student
recognizability. Limited. A real bilateral trade panel covers 200
territories and goes back decades. Whatever your slice is, the scope is
part of every claim you can make.

\textbf{6. Where did the data come from, and what survived the journey?}
Raw data is messy. Clean data is the result of decisions: which
observations to drop, which variables to re-code, which missing values
to impute, which outliers to keep or remove. Every one of those
decisions is a place where signal can leak. The processing script is
part of the data. So is the person who wrote it.

\textbf{7. What is missing?} The thing the dataset doesn't have is often
more important than the thing it does. Our slice has six countries.
Missing: every country except those six, every year before 2014, every
product category, every direction of investment. Any conclusion drawn
from this slice is conditional on those exclusions. A good analyst names
the absences before drawing the inferences.

\textbf{8. What looks weird?} Look at the data. Sort by exports, find
the smallest and largest values, check whether they make sense. Look at
the distance column, find the row where two countries are improbably
close or far, ask why. Look at the GDP columns, look for off-by-1000
unit confusion. The first ten minutes of weirdness-hunting saves the
next ten hours.

These eight questions will not catch every problem. They will catch
most. Run them every time. The discipline is what separates analysis
from rationalization.

\begin{tcolorbox}[enhanced jigsaw, arc=.35mm, bottomrule=.15mm, bottomtitle=1mm, breakable, colback=white, colbacktitle=quarto-callout-note-color!10!white, colframe=quarto-callout-note-color-frame, coltitle=black, left=2mm, leftrule=.75mm, opacityback=0, opacitybacktitle=0.6, rightrule=.15mm, title=\textcolor{quarto-callout-note-color}{\faInfo}\hspace{0.5em}{Inside the math: tidy data, formally}, titlerule=0mm, toprule=.15mm, toptitle=1mm]

The rectangular structure described in this chapter is one specific
arrangement called \emph{tidy data}, formalized by Wickham (2014). A
dataset is tidy when:

\begin{enumerate}
\def\labelenumi{\arabic{enumi}.}
\tightlist
\item
  Each variable forms a column.
\item
  Each observation forms a row.
\item
  Each type of observational unit forms a table.
\end{enumerate}

These three conditions are equivalent to a special kind of relational
normal form. The first two are the rectangular structure we have been
discussing. The third is the rule that distinct kinds of things
(countries, country-pairs, years) live in distinct tables, joined by
keys, rather than crammed into a single table with mixed grain.

Many real datasets are not tidy. Wide-format panels, wide-format
surveys, and aggregated tables encountered in the wild often fail one or
more conditions. Most data analysis begins with tidying: reshaping data
into the rectangular form. Most data analysis errors begin with someone
forgetting which form they are in.

The reference is Hadley Wickham, ``Tidy Data,'' \emph{Journal of
Statistical Software}, 59(10), 2014.

\end{tcolorbox}

\begin{tcolorbox}[enhanced jigsaw, arc=.35mm, bottomrule=.15mm, bottomtitle=1mm, breakable, colback=white, colbacktitle=quarto-callout-tip-color!10!white, colframe=quarto-callout-tip-color-frame, coltitle=black, left=2mm, leftrule=.75mm, opacityback=0, opacitybacktitle=0.6, rightrule=.15mm, title=\textcolor{quarto-callout-tip-color}{\faLightbulb}\hspace{0.5em}{Try this}, titlerule=0mm, toprule=.15mm, toptitle=1mm]

Open \texttt{data/trade/clean/bilateral\_panel.csv} in any tool you like
(a spreadsheet, R, Python, or your text editor of choice; the file is
small).

Run the eight-question audit on it. Write your answers down. The point
is not to be exhaustive; the point is to make the questions habitual.
When you finish, two things should be true: you can describe the dataset
to someone in three sentences, and you can list at least one thing you
wish were in there but is not.

If you want a stretch question: compute the average bilateral exports
for each year, plot the result, and notice what happens around 2020.
Speculate about why.

\end{tcolorbox}

\begin{tcolorbox}[enhanced jigsaw, arc=.35mm, bottomrule=.15mm, bottomtitle=1mm, breakable, colback=white, colbacktitle=quarto-callout-important-color!10!white, colframe=quarto-callout-important-color-frame, coltitle=black, left=2mm, leftrule=.75mm, opacityback=0, opacitybacktitle=0.6, rightrule=.15mm, title=\textcolor{quarto-callout-important-color}{\faExclamation}\hspace{0.5em}{Frontier note}, titlerule=0mm, toprule=.15mm, toptitle=1mm]

The discipline of asking what data is and where it came from has a
formal name in modern empirical work: \emph{data provenance}. A growing
literature in measurement theory and metascience argues that most
published quantitative claims are weaker than they appear because the
chain from raw observation to reported number is opaque, and that
openness about that chain (raw data, processing scripts, documented
decisions) should be a precondition for citation.

Two places to follow this if you care: the Berkeley Initiative for
Transparency in the Social Sciences (BITSS), and the literature growing
out of the credibility revolution in economics. Both have produced
practical tools and standards for the kind of provenance work this
chapter has gestured toward.

\end{tcolorbox}

\begin{tcolorbox}[enhanced jigsaw, arc=.35mm, bottomrule=.15mm, bottomtitle=1mm, breakable, colback=white, colbacktitle=quarto-callout-warning-color!10!white, colframe=quarto-callout-warning-color-frame, coltitle=black, left=2mm, leftrule=.75mm, opacityback=0, opacitybacktitle=0.6, rightrule=.15mm, title=\textcolor{quarto-callout-warning-color}{\faExclamationTriangle}\hspace{0.5em}{Key insight}, titlerule=0mm, toprule=.15mm, toptitle=1mm]

Data is rectangular. Rows are observations, columns are variables, and
the unit of analysis is the question that has to be answered before
anything else. The eight-question audit is the discipline that turns a
stack of numbers into evidence.

\end{tcolorbox}

\section{Practice problems}\label{practice-problems}

\begin{enumerate}
\def\labelenumi{\arabic{enumi}.}
\item
  Take any spreadsheet you can put your hands on. Run the eight-question
  audit on it. Write your answers in three sentences each. The
  discipline gets faster after about three datasets.
\item
  Our trade slice has unit of analysis (exporter, importer, year).
  Suggest two other reasonable units of analysis you could aggregate to.
  For each, name one question that becomes easier and one that becomes
  impossible.
\item
  A column called \texttt{blood\_type} has values A+, A−, B+, B−, AB+,
  AB−, O+, O−. What variable type is it? Could you ever defend treating
  it as ordinal?
\item
  ZIP codes look like integers, and computing the mean ZIP code of a
  sample is mathematically valid. Explain in two sentences why the
  result is useless.
\item
  Open \texttt{data/trade/codebook.md}. Identify two facts about the
  dataset that you would not have learned by looking at the CSV alone.
\end{enumerate}

\begin{tcolorbox}[enhanced jigsaw, arc=.35mm, bottomrule=.15mm, bottomtitle=1mm, breakable, colback=white, colbacktitle=quarto-callout-caution-color!10!white, colframe=quarto-callout-caution-color-frame, coltitle=black, left=2mm, leftrule=.75mm, opacityback=0, opacitybacktitle=0.6, rightrule=.15mm, title=\textcolor{quarto-callout-caution-color}{\faFire}\hspace{0.5em}{Reflect before moving on}, titlerule=0mm, toprule=.15mm, toptitle=1mm]

\begin{itemize}
\tightlist
\item
  In your own words, what is the difference between ``rows are
  observations'' and ``rows are independent''?
\item
  Why might the same dataset be tidy for one analysis and not tidy for
  another?
\item
  Pick the audit question you would be most likely to skip in practice.
  Why? What do you risk by skipping it?
\end{itemize}

\end{tcolorbox}

\chapter{Probability as area}\label{probability-as-area}

\section{What probability is}\label{what-probability-is}

Pause on the word ``probability.'' Pretend you have never heard it
before. What kind of thing is it?

Here is the answer the book will use, and it is the answer modern
probability theory uses too: \emph{probability is a fraction of an
outcome space.}

That sentence does a lot of work. Let me unpack it slowly.

An \emph{outcome space} is the set of all the things that could happen
in some situation. When you roll a fair die, the outcome space is the
six numbers from 1 to 6. When you draw a card from a shuffled deck, the
outcome space is the 52 distinct cards. When you flip a coin, the
outcome space has two elements. When you measure tomorrow's high
temperature in Atlanta, the outcome space is technically every real
number between some lower and upper bound, an infinite continuum.

Once you have an outcome space, probability is just bookkeeping over the
space. The probability of any specific outcome is a number between 0 and
1. Add up the probabilities of the outcomes inside any subset, and you
get that subset's probability. The probability of the \emph{whole space}
is exactly 1, because something has to happen.

That is it. There is a deeper formal version with σ-algebras and measure
theory waiting later in the book, but at this level, probability is
``what fraction of the things that could happen does this subset
cover.''

The picture below, and the rest of this chapter, takes that definition
seriously.

\begin{figure*}

\centering{

\pandocbounded{\includegraphics[keepaspectratio]{prose/part1/../../pictures/figures/output/fig_sample_space.pdf}}

}

\caption{\label{fig-sample-space}Every outcome of rolling two fair dice.
There are 36 cells, one for each ordered pair of die-1 and die-2 values,
and the number in each cell is the sum. The shading shows how many cells
share the same sum: the diagonal stripe with sum 7 has six cells, sum 6
and sum 8 have five, and so on, down to sums 2 and 12 with one cell
each. Probability becomes geometry: the chance of rolling a 7 is just
the fraction of the rectangle that is shaded most darkly.}

\end{figure*}%

Notice what the picture does not yet have: any equation. Probability is
not, fundamentally, an equation. It is a way of carving up a space. The
equations come later, as shorthand for the carving.

The same idea works whether the space is finite or continuous. When you
throw a dart at a square dartboard, the outcome space is the dartboard
itself. The probability that the dart lands in some region is the
fraction of the dartboard that region covers. This is literally area
divided by area, no abstraction required.

When you measure tomorrow's temperature, the outcome space is a number
line and the probabilities are encoded by a \emph{density}: a curve over
the number line whose total area equals 1. The probability that the
temperature lands in some interval is the area under the density curve
over that interval. We will spend chapter three meeting the most
important density curves and learning to read them. For now, the only
thing to remember is that area under a curve is probability. Always. No
exceptions.

\section{Events are regions}\label{events-are-regions}

An \emph{event} is just a name for a subset of the outcome space.

``The card I draw is a heart'' is an event. It corresponds to the subset
of the deck consisting of the 13 hearts. The probability of the event is
13/52, or 1/4, because the heart-region covers a quarter of the deck.

``The card is a face card'' (jack, queen, king of any suit) is another
event. It corresponds to 12 cards out of 52. Probability 12/52, about
0.23.

What about ``the card is a heart \textbf{and} a face card''? That is the
\emph{intersection} of the two events: the cards that are simultaneously
hearts and face cards. There are exactly three of those (jack, queen,
king of hearts). Probability 3/52, about 0.058.

What about ``the card is a heart \textbf{or} a face card''? That is the
\emph{union}. You might be tempted to add the two probabilities, 13/52 +
12/52 = 25/52. That is wrong, and the reason it is wrong is one of the
most useful diagrams in elementary probability.

The hearts and the face cards overlap. The three cards in the
intersection (J, Q, K of hearts) belong to both regions. If you add 13
and 12, you have counted those three twice. To get the union, you add
the two regions and subtract the overlap once:

\[
P(\text{heart} \cup \text{face}) = P(\text{heart}) + P(\text{face}) - P(\text{heart} \cap \text{face}) = \frac{13}{52} + \frac{12}{52} - \frac{3}{52} = \frac{22}{52}.
\]

This is the \emph{inclusion-exclusion principle}, and it has a name only
because beginners get it wrong so often. The picture is unambiguous.

\begin{figure*}

\centering{

\pandocbounded{\includegraphics[keepaspectratio]{prose/part1/../../pictures/figures/output/fig_events_overlap.pdf}}

}

\caption{\label{fig-events-overlap}Two events drawn as overlapping
regions inside a rectangular sample space. The blue region A is ``the
card is a heart'' (13 of 52). The rust region B is ``the card is a face
card'' (12 of 52). The overlap A ∩ B has three cards (J, Q, K of
hearts). When you compute P(A ∪ B), the overlap is counted once in P(A)
and once in P(B), which means you must subtract it once to avoid
double-counting.}

\end{figure*}%

Two events that don't overlap are called \emph{disjoint}, or
\emph{mutually exclusive}. For disjoint events, the union really is just
the sum, because there is nothing to subtract. Drawing a heart and
drawing a club are disjoint: no card is both. P(heart or club) = 13/52 +
13/52 = 26/52, exactly half the deck.

Read the figure carefully. Every claim in the next two chapters can be
derived from looking at it.

\section{Conditional probability is a change of
denominator}\label{conditional-probability-is-a-change-of-denominator}

Suppose you draw a card and someone tells you, before you look, that the
card is a heart. Now what is the probability the card is a face card?

The original outcome space was 52 cards, of which 12 were face cards.
P(face) = 12/52, about 0.23. But you have new information: the card is a
heart. The outcome space has shrunk. Instead of 52 cards, only the 13
hearts are now possibilities. Among those 13, exactly 3 are face cards
(J, Q, K of hearts). The probability that the card is a face card
\emph{given that it is a heart} is 3/13, about 0.23.

Conditioning on a piece of information is just zooming in. You take the
original outcome space, restrict it to the region you have been told
about, and recompute the fraction.

\begin{figure*}

\centering{

\pandocbounded{\includegraphics[keepaspectratio]{prose/part1/../../pictures/figures/output/fig_conditional_zoom.pdf}}

}

\caption{\label{fig-cond-zoom}Left: the full deck, with the 12 face
cards highlighted in gold. The probability of a face card, P(face), is
12/52 ≈ 0.23. Right: the same problem, restricted to the 13 hearts only.
The probability of a face card given that the card is a heart, P(face
given heart), is 3/13 ≈ 0.23. Conditioning on ``heart'' zoomed the
sample space from 52 cards to 13. The numerator and denominator both
changed; in this case the ratio happens to come out about the same.}

\end{figure*}%

Notation. The probability of A given B is written \(P(A \mid B)\), read
aloud as ``the probability of A given B.'' The vertical bar means
``given.'' It is the most important piece of notation in inferential
statistics, and most students never learn to read it aloud confidently.
Read it aloud now.

The formula corresponding to the picture is

\[
P(A \mid B) = \frac{P(A \cap B)}{P(B)}.
\]

Read the formula aloud. \emph{The probability of A given B equals the
probability of A and B together, divided by the probability of B.} In
words: of all the times B happened, what fraction of them also had A
happen?

Verify it on the cards.
\(P(\text{face} \mid \text{heart}) = P(\text{face} \cap \text{heart}) / P(\text{heart}) = (3/52) / (13/52) = 3/13\).
Same answer the picture gave. The formula and the picture are saying the
same thing in different languages.

Internalize this: \(P(A \mid B)\) is a different quantity from \(P(A)\).
The two probabilities can be very different. Sometimes B raises the
chance of A; sometimes B lowers it. When the two probabilities happen to
be equal, A and B are \emph{independent}, and B carries no information
about A. The same word ``probability'' is doing very different work in
\(P(A)\) and \(P(A \mid B)\), and a great deal of statistical confusion
in the wild comes from people sliding between them without noticing.

\section{Bayes and the disease-test
paradox}\label{bayes-and-the-disease-test-paradox}

Here is a question that broke me the first time I saw it, and that I
have watched humble dozens of trained scientists since.

A rare disease affects 1\% of the population. A test for the disease is
95\% accurate: among people who have the disease, the test correctly
flags 95\% of them; among people who do not have the disease, the test
correctly clears 95\% of them. You take the test. It comes back
positive.

What is the probability you have the disease?

Most people, when surveyed, answer somewhere between 90\% and 95\%,
anchoring on the test's accuracy. The actual answer, as the picture
below makes obvious, is roughly 16\%.

Imagine 10,000 people drawn at random from the population. The picture
shows them as a 100-by-100 grid of squares.

\begin{figure*}

\centering{

\pandocbounded{\includegraphics[keepaspectratio]{prose/part1/../../pictures/figures/output/fig_bayes_grid.pdf}}

}

\caption{\label{fig-bayes-grid}10,000 people, drawn at random from a
population with 1\% disease prevalence and tested with a 95\%-accurate
test. The dark red squares are sick people who tested positive (95). The
pale red squares are sick people who tested negative (5). The rust
squares are healthy people who tested positive anyway, the false
positives (495). The gray squares are healthy people who tested negative
(9,405). The probability of being sick given a positive test is the
fraction of the positive-test population that is also in the sick
population: 95 out of (95 + 495) = 590, or about 16\%. The base rate
dominates.}

\end{figure*}%

Of the 10,000 people, 1\% have the disease. That is 100 people. The
other 9,900 do not.

Among the 100 sick people, the test correctly identifies 95\% as
positive. That is 95 true positives. The other 5 are false negatives.

Among the 9,900 well people, the test correctly clears 95\%. That is
9,405 true negatives. But it incorrectly flags 5\%. That is 495 false
positives.

Now your test is positive. You belong to the population that tests
positive. That population has two parts: the 95 true positives (sick)
and the 495 false positives (well). The whole positive-test population
is 590 people.

Of those 590, only 95 are actually sick. The probability that a
positive-test person has the disease is 95/590, about 0.161, about 16\%.

The 95\% accuracy did not lie. It just answered a different question.
\(P(+ \mid \text{sick}) = 0.95\) is not the same number as
\(P(\text{sick} \mid +)\). The base rate, that 99\% of the population is
well, dominates the calculation. With a rare disease and a test that is
anything less than perfect, most positives are false positives.

This is \emph{Bayes' theorem}, and the formula version is

\[
P(\text{sick} \mid +) = \frac{P(+ \mid \text{sick}) \, P(\text{sick})}{P(+)} = \frac{0.95 \times 0.01}{0.95 \times 0.01 + 0.05 \times 0.99} \approx 0.161.
\]

Read it aloud. \emph{The probability of being sick given a positive test
equals the probability of testing positive given sickness, times the
prevalence of sickness, divided by the total probability of testing
positive.} That denominator is just the size of the positive-test
population (the 590 from the picture), expressed as a fraction.

The formula and the picture say exactly the same thing. The picture is
more honest. When you read about Bayes' theorem in any later chapter,
picture the grid first and translate the formula into rectangles second.
That is the pedagogical move that prevents an entire class of mistakes
that even careful researchers continue to make.

\begin{tcolorbox}[enhanced jigsaw, arc=.35mm, bottomrule=.15mm, bottomtitle=1mm, breakable, colback=white, colbacktitle=quarto-callout-warning-color!10!white, colframe=quarto-callout-warning-color-frame, coltitle=black, left=2mm, leftrule=.75mm, opacityback=0, opacitybacktitle=0.6, rightrule=.15mm, title=\textcolor{quarto-callout-warning-color}{\faExclamationTriangle}\hspace{0.5em}{Key insight}, titlerule=0mm, toprule=.15mm, toptitle=1mm]

Probability is a fraction of an outcome space. Conditioning is
restricting that space. Bayes' theorem is the bookkeeping that connects
the two, and the rectangle-grid picture in this chapter is the only
intuition you need to apply it correctly in the simple case.

\end{tcolorbox}

\begin{tcolorbox}[enhanced jigsaw, arc=.35mm, bottomrule=.15mm, bottomtitle=1mm, breakable, colback=white, colbacktitle=quarto-callout-tip-color!10!white, colframe=quarto-callout-tip-color-frame, coltitle=black, left=2mm, leftrule=.75mm, opacityback=0, opacitybacktitle=0.6, rightrule=.15mm, title=\textcolor{quarto-callout-tip-color}{\faLightbulb}\hspace{0.5em}{Try this}, titlerule=0mm, toprule=.15mm, toptitle=1mm]

A different version of the disease test problem. The disease prevalence
is 50\% (very common). The test accuracy is the same 95\%.

Without computing anything, predict whether \(P(\text{sick} \mid +)\) is
closer to 50\% or 95\%.

Then draw the grid. Then compute the answer. Compare to your prediction.
Notice how much the prevalence is doing in this problem, even though the
prevalence never appears in the test's specifications.

A second one. Same 1\% prevalence as before, but now imagine the test is
99\% accurate (99\% sensitivity, 99\% specificity). Recompute
\(P(\text{sick} \mid +)\). What changed and what didn't?

\end{tcolorbox}

\begin{tcolorbox}[enhanced jigsaw, arc=.35mm, bottomrule=.15mm, bottomtitle=1mm, breakable, colback=white, colbacktitle=quarto-callout-note-color!10!white, colframe=quarto-callout-note-color-frame, coltitle=black, left=2mm, leftrule=.75mm, opacityback=0, opacitybacktitle=0.6, rightrule=.15mm, title=\textcolor{quarto-callout-note-color}{\faInfo}\hspace{0.5em}{Inside the math: Bayes from first principles}, titlerule=0mm, toprule=.15mm, toptitle=1mm]

The chapter derived Bayes' theorem from a picture. Here is the same
derivation in symbols, for readers who want to see why it always works.

The conditional probability formula \(P(A \mid B) = P(A \cap B) / P(B)\)
can be rearranged into \(P(A \cap B) = P(A \mid B) \, P(B)\). This is
sometimes called the \emph{multiplication rule}. Notice that the same
intersection \(A \cap B\) can be written by conditioning the other
direction: \(P(A \cap B) = P(B \mid A) \, P(A)\). Setting the two equal,

\[
P(A \mid B) \, P(B) = P(B \mid A) \, P(A),
\]

and dividing both sides by \(P(B)\) gives Bayes' theorem in its simplest
form,

\[
P(A \mid B) = \frac{P(B \mid A) \, P(A)}{P(B)}.
\]

The denominator \(P(B)\) is the \emph{marginal probability} of B, and in
problems like the disease test it is computed by summing over the
conditioning variable:
\(P(B) = P(B \mid A) P(A) + P(B \mid \text{not } A) P(\text{not } A)\).
This is the \emph{law of total probability}, and it is the formal
version of ``add up the squares in the grid that meet condition B.''

That is the entire mathematical content of Bayesian inference.
Everything that follows in chapters 30 and beyond is just turning the
rectangles into curves and the counts into integrals.

\end{tcolorbox}

\begin{tcolorbox}[enhanced jigsaw, arc=.35mm, bottomrule=.15mm, bottomtitle=1mm, breakable, colback=white, colbacktitle=quarto-callout-important-color!10!white, colframe=quarto-callout-important-color-frame, coltitle=black, left=2mm, leftrule=.75mm, opacityback=0, opacitybacktitle=0.6, rightrule=.15mm, title=\textcolor{quarto-callout-important-color}{\faExclamation}\hspace{0.5em}{Frontier note}, titlerule=0mm, toprule=.15mm, toptitle=1mm]

Modern Bayesian statistics is the largest practical descendant of the
simple bookkeeping you just learned. A working Bayesian analysis turns
the rectangle picture from this chapter into a continuous version (the
rectangles become curves and the counts become integrals), and asks the
same question: given the data, what is the conditional distribution of
the parameters?

The recent renaissance of Bayesian methods in applied work owes most of
its momentum to two things: better computational tools (Hamiltonian
Monte Carlo, variational inference), and a slow shift in scientific
culture toward thinking about uncertainty in the way Bayes' theorem
requires.

For now, every time you condition on data in this book, remember that
you are doing a tiny piece of Bayesian thinking, even if the chapter
calls it something else.

\end{tcolorbox}

\section{Practice problems}\label{practice-problems-1}

\begin{enumerate}
\def\labelenumi{\arabic{enumi}.}
\item
  A bag has 5 red, 3 blue, and 2 green marbles. You draw one. Compute
  \(P(\text{red})\), \(P(\text{blue or green})\), and
  \(P(\text{red and blue})\). Sketch the sample space as a rectangle to
  verify.
\item
  Repeat the disease-test problem with prevalence 5\% and test accuracy
  90\% (90\% sensitivity, 90\% specificity). Draw the grid for 10,000
  people and compute \(P(\text{sick} \mid +)\).
\item
  Two events \(A\) and \(B\) have \(P(A) = 0.4\), \(P(B) = 0.5\),
  \(P(A \cap B) = 0.1\). Compute \(P(A \cup B)\), \(P(A \mid B)\), and
  \(P(B \mid A)\). Are \(A\) and \(B\) independent?
\item
  Why is \(P(A \cup B) = P(A) + P(B) - P(A \cap B)\) rather than just
  \(P(A) + P(B)\)? Explain in one sentence and a small picture.
\item
  A test has 99\% sensitivity and 99\% specificity, and the underlying
  disease has prevalence 0.1\%. Without computing, predict whether
  \(P(\text{sick} \mid +)\) is closer to 99\% or 10\%. Then compute and
  compare.
\end{enumerate}

\begin{tcolorbox}[enhanced jigsaw, arc=.35mm, bottomrule=.15mm, bottomtitle=1mm, breakable, colback=white, colbacktitle=quarto-callout-caution-color!10!white, colframe=quarto-callout-caution-color-frame, coltitle=black, left=2mm, leftrule=.75mm, opacityback=0, opacitybacktitle=0.6, rightrule=.15mm, title=\textcolor{quarto-callout-caution-color}{\faFire}\hspace{0.5em}{Reflect before moving on}, titlerule=0mm, toprule=.15mm, toptitle=1mm]

\begin{itemize}
\tightlist
\item
  What is the visual difference between conditional probability and
  joint probability?
\item
  Where else in the world have you seen base-rate neglect at work?
\item
  If a friend asked you ``what is probability, in one sentence,'' what
  would you say now that you could not have said before this chapter?
\end{itemize}

\end{tcolorbox}

\chapter{Shapes of data}\label{shapes-of-data}

\section{What a distribution is}\label{what-a-distribution-is}

You have a column. A few hundred rows of bilateral exports, or trade
distances, or human heights, or wait times for the bus. The values are
not all the same, and they are not arranged at random across some flat
range. They cluster, they spread, they trail off. The pattern of where
the values fall is what statisticians call the \emph{distribution}.

A distribution is the answer to ``for any given value, how often does it
appear?'' If you collect heights of adult men in centimeters and bin
them into one-centimeter buckets, you can plot the count in each bucket
and watch a shape emerge. The shape is not random. It has a name. It
carries information about the kind of process that produced the values.

Most of the conceptual work of statistics is recognizing distribution
shapes and matching them to processes. You learn the half-dozen most
important shapes and which kinds of real-world situations produce them,
and a great deal of the rest of the field becomes the question of how to
use those shapes to make good decisions.

This chapter introduces six distributions you will see repeatedly. None
of them are formulas to memorize. Each is a \emph{story}: a description
of what generates data with this shape. When you meet a column of
numbers in the wild, your first move is to ask what story produced it.
The right answer points at the right tool.

Six stories follow. Each gets a figure. Read the figures first, the
prose second.

\section{Bernoulli and binomial: a coin and a
tally}\label{bernoulli-and-binomial-a-coin-and-a-tally}

Start with the simplest possible distribution. You have a single
yes-or-no event with probability \(p\) of being yes. Coin flip with a
biased coin. Click-or-no-click on an ad. Whether a randomly chosen
person owns a passport. The outcome is one of two values, conventionally
encoded as 1 (yes) and 0 (no). The distribution is called
\emph{Bernoulli}, and it is fully described by the single number \(p\).

You can plot a Bernoulli distribution as two bars: one of height \(p\)
at the value 1, one of height \(1-p\) at the value 0. That is the entire
shape. There is nothing more to know.

Now suppose you do \(n\) independent Bernoulli trials with the same
\(p\) and count how many came out yes. The number of yeses is no longer
a yes-or-no; it is an integer between 0 and \(n\). The distribution of
that count is called \emph{binomial}, written Binomial(\(n, p\)). It has
more shape than the Bernoulli does, because the count can land anywhere
in a range, and some counts are far more likely than others.

When \(n\) is small, the binomial is asymmetric and lumpy. When \(n\) is
large, it starts to look like a bell curve centered at \(np\) with width
controlled by \(\sqrt{np(1-p)}\). We will return to that ``starts to
look like'' later in the book, because it is one of the most important
facts in statistics. For now, recognize the shape.

\begin{figure*}

\centering{

\pandocbounded{\includegraphics[keepaspectratio]{prose/part1/../../pictures/figures/output/fig_bernoulli_binomial.pdf}}

}

\caption{\label{fig-bernoulli-binomial}Top: a Bernoulli distribution at
three values of p (0.2, 0.5, 0.8). It is just two bars, one at 0 and one
at 1. Bottom: a Binomial(20, p) at the same three values. Notice how the
binomial spreads, peaks near np, and is symmetric only when p = 0.5.}

\end{figure*}%

Where you meet binomials in the wild: anything counted from a fixed
number of trials. Test scores out of 50. Free throws made in 10
attempts. People in a sample of 1,000 who report a given opinion. The
story is always ``fixed \(n\), fixed \(p\), count the yeses.''

\section{Poisson: the rare event}\label{poisson-the-rare-event}

A different story. You are not running a fixed number of trials; you are
watching a clock or a region in space, and events happen at some average
rate. Earthquakes per year above some magnitude. Cars arriving at a toll
booth per minute. Typos per page. Cosmic rays striking your phone in an
hour. Events are independent, the rate is approximately constant, and
you count how many you observed in your window.

The distribution of that count is called \emph{Poisson}, written
Poisson(\(\lambda\)). Its single parameter \(\lambda\) is the average
number of events expected in the window. Values live on the non-negative
integers (0, 1, 2, \ldots), with a characteristic shape that peaks near
\(\lambda\), drops off on both sides, and never quite reaches zero.

\begin{figure}

\centering{

\pandocbounded{\includegraphics[keepaspectratio]{prose/part1/../../pictures/figures/output/fig_poisson.pdf}}

}

\caption{\label{fig-poisson}Poisson PMFs at three values of \(\lambda\).
Small \(\lambda\) produces a heavily right-skewed shape with most of the
mass at zero or one. As \(\lambda\) grows, the distribution shifts right
and starts to look more like a bell.}

\end{figure}%

The Poisson story is the answer to ``how many rare independent events in
a fixed window?'' When you see count data and the events are not all
going to be yes-or-no with a fixed \(n\), the Poisson is usually the
right starting point. Many real count distributions are not exactly
Poisson but are close enough that the Poisson is the right idealization.

A useful trick. If \(\lambda\) is large (say, more than 30), the Poisson
starts to look like a normal distribution centered at \(\lambda\) with
standard deviation \(\sqrt{\lambda}\). This is a special case of a
phenomenon we will spend chapter 14 working on.

\section{Normal: the sum of many small
things}\label{normal-the-sum-of-many-small-things}

The most important distribution in statistics, full stop. The bell
curve. The Gaussian. The normal. There are reasons it shows up
everywhere, and the reasons matter more than the formula.

The story: a normal distribution arises whenever the value you are
measuring is the sum of many independent small contributions, none of
which dominates the total. Human height is approximately normal because
it is the sum of many small genetic and environmental influences.
Measurement error is approximately normal because each measurement is a
sum of many tiny independent perturbations. The mean of any sample (no
matter what the underlying distribution is) is approximately normal once
\(n\) is large enough. We will spend chapter 14 on that last fact,
because it is the engine of inferential statistics.

A normal distribution is fully described by two numbers: its mean (where
it is centered) and its standard deviation (how wide it is). The shape
is symmetric about the mean, drops off rapidly on both sides, and has
the characteristic 68-95-99.7 property: about 68\% of the area is within
one standard deviation of the mean, about 95\% within two, about 99.7\%
within three.

\begin{figure}

\centering{

\pandocbounded{\includegraphics[keepaspectratio]{prose/part1/../../pictures/figures/output/fig_normal_68_95_99.pdf}}

}

\caption{\label{fig-normal}A standard normal distribution with the
68-95-99.7 bands shaded. Reading the figure: roughly two-thirds of the
area falls inside the inner band, ninety-five percent inside the middle
band, and almost all the rest inside the outer band. The thresholds at
\(\pm 2\sigma\) and \(\pm 3\sigma\) are the natural reference points for
any normal-shaped data.}

\end{figure}%

The 68-95-99.7 rule is worth memorizing. It lets you do quick sanity
checks on any normal-shaped data: if a value is more than 2 sd from the
mean, it is in the top or bottom 2.5\% of the distribution. More than 3
sd, top or bottom 0.15\%. The thresholds become natural after a few
uses.

\section{Exponential and log-normal: the wait, the
multiplied}\label{exponential-and-log-normal-the-wait-the-multiplied}

Two more shapes that recur, both heavy on the right side.

The \emph{exponential} distribution is the story of waiting. You are
watching for the next earthquake, the next bus, the next user click.
Time elapses. The probability of waiting at least \(t\) more time units
shrinks at a constant rate, regardless of how long you have already
waited (the \emph{memoryless} property). The distribution of the wait
time has a peak at zero and a fast decay toward longer waits. Most waits
are short, a few are surprisingly long. The single parameter is the rate
\(\lambda\), and the mean wait is \(1/\lambda\).

The \emph{log-normal} distribution is the story of multiplied effects.
Suppose a quantity is built up not by adding many small things but by
multiplying them. The size of a city is roughly the product of
historical advantages, each one acting as a multiplier. Income is
roughly the product of skill, opportunity, market timing, and luck, each
multiplying the previous. The log of such a quantity is a sum of small
things, which (by the previous section's story) is approximately normal.
So the original quantity is the exponential of a normal, called
\emph{log-normal}.

\begin{figure*}

\centering{

\pandocbounded{\includegraphics[keepaspectratio]{prose/part1/../../pictures/figures/output/fig_exp_lognormal.pdf}}

}

\caption{\label{fig-exp-lognormal}Left: an exponential distribution.
Most of the mass is at small values; a thin tail stretches to the right.
Right: a log-normal distribution. Heavier on the right than the
exponential; a small fraction of values are extreme. Income, city
populations, file sizes, reaction times in psychology: all approximately
log-normal.}

\end{figure*}%

The first move when you see right-skew data is usually to take the log
and see if the result looks normal. Often it does. Once it does, all the
tools you have for working with normal distributions become available,
applied on the log scale. Coefficients of regressions on log-transformed
outcomes have a natural interpretation as percentage changes, which we
will use heavily in the regression chapters.

\section{The t: when n is small}\label{the-t-when-n-is-small}

One last shape. The \emph{Student's t} distribution, which looks almost
like a normal but is not.

The story: when you compute the mean of a small sample, you do not
actually know the population standard deviation; you have to estimate it
from the data. That estimate has its own uncertainty, especially when
the sample is small. The t distribution is the correct accounting for
the extra uncertainty. It looks like a normal but with heavier tails,
the heaviness controlled by a parameter called \emph{degrees of freedom}
that essentially equals the sample size minus one.

\begin{figure}

\centering{

\pandocbounded{\includegraphics[keepaspectratio]{prose/part1/../../pictures/figures/output/fig_t_vs_normal.pdf}}

}

\caption{\label{fig-t-vs-normal}The t distribution overlaid on a
standard normal. At \(df = 2\), the t is dramatically heavier-tailed
than the normal: substantially more mass beyond \(\pm 2\). At
\(df = 10\) the difference is smaller but still visible. By \(df = 30\)
the t is almost indistinguishable from the normal.}

\end{figure}%

When the sample is large (say, \(n > 30\)), the t is indistinguishable
from the normal. When the sample is small (say, \(n = 5\)), the t has
noticeably fatter tails: more probability of extreme values, to account
for the fact that the standard deviation estimate could itself be off.
As \(n\) grows, the t converges to the normal.

You will meet the t the first time we do a confidence interval for a
mean, in chapter 13. The reason hypothesis tests for small samples use
the t and not the normal is exactly the story above.

\begin{tcolorbox}[enhanced jigsaw, arc=.35mm, bottomrule=.15mm, bottomtitle=1mm, breakable, colback=white, colbacktitle=quarto-callout-warning-color!10!white, colframe=quarto-callout-warning-color-frame, coltitle=black, left=2mm, leftrule=.75mm, opacityback=0, opacitybacktitle=0.6, rightrule=.15mm, title=\textcolor{quarto-callout-warning-color}{\faExclamationTriangle}\hspace{0.5em}{Key insight}, titlerule=0mm, toprule=.15mm, toptitle=1mm]

Distributions are stories about generating processes. Six of them cover
most of what you will see in working analysis: Bernoulli (yes-or-no),
Binomial (counts of yeses), Poisson (rare events in a window), Normal
(sums of many small things), Exponential and Log-normal (waits and
products), and Student's t (the normal's small-sample cousin).
Recognizing the story tells you which tool to use; the formulas come
later.

\end{tcolorbox}

\begin{tcolorbox}[enhanced jigsaw, arc=.35mm, bottomrule=.15mm, bottomtitle=1mm, breakable, colback=white, colbacktitle=quarto-callout-tip-color!10!white, colframe=quarto-callout-tip-color-frame, coltitle=black, left=2mm, leftrule=.75mm, opacityback=0, opacitybacktitle=0.6, rightrule=.15mm, title=\textcolor{quarto-callout-tip-color}{\faLightbulb}\hspace{0.5em}{Try this}, titlerule=0mm, toprule=.15mm, toptitle=1mm]

Pick a column from any dataset you have access to, ideally something
with at least a hundred values. Plot a histogram. Ask which of the six
stories above best describes its shape. Then ask why: what process
produced the values, and does that process match the story?

If the data is right-skewed, take the log and plot again. Does the
log-transformed version look approximately normal? If yes, you have a
log-normal column on your hands and a different set of tools applies.

A second prompt: pick two columns from the trade slice
(\texttt{exports\_usd\_millions} and \texttt{distance\_km} are good).
Plot a histogram of each. Then plot a histogram of
\(\log(\text{exports})\) and \(\log(\text{distance})\). Notice which
transformation does meaningful work.

\end{tcolorbox}

\begin{tcolorbox}[enhanced jigsaw, arc=.35mm, bottomrule=.15mm, bottomtitle=1mm, breakable, colback=white, colbacktitle=quarto-callout-note-color!10!white, colframe=quarto-callout-note-color-frame, coltitle=black, left=2mm, leftrule=.75mm, opacityback=0, opacitybacktitle=0.6, rightrule=.15mm, title=\textcolor{quarto-callout-note-color}{\faInfo}\hspace{0.5em}{Inside the math: where the names come from}, titlerule=0mm, toprule=.15mm, toptitle=1mm]

Each distribution in this chapter has a probability mass function (for
discrete distributions) or probability density function (for continuous
ones). The names attached to the formulas are sedimented history.

Jacob Bernoulli, working on games of chance in the late 1600s, gave his
name to the simplest yes-or-no distribution. The binomial coefficient
that appears inside Binomial(\(n, p\)) was studied earlier still, by
Pascal and Fermat in their 1654 correspondence about gambling.

In 1837, Siméon Denis Poisson published a treatise on the probability of
judicial verdicts that contained what we now call the \emph{Poisson
distribution}. Its earliest serious application was to the statistics of
rare events in human populations: deaths from horse-kicks in the
Prussian army, used by Bortkiewicz in 1898 as a textbook example.

The normal distribution was discovered three times. De Moivre derived it
in 1733 as a limit of the binomial. Laplace developed it independently
in 1810 inside what we now call the central limit theorem. Gauss arrived
at it from a different direction in 1809, in his theory of measurement
error. ``Normal'' as the popular name came only later, from Galton and
Pearson at the end of the 1800s.

In 1908, a Guinness brewery statistician named William Sealy Gosset
published the t distribution under the pen name \emph{Student}, because
his employer forbade publication under his real name. It is one of the
most consequential pseudonyms in scientific history.

Log-normal and exponential are simpler in provenance: both fall out
naturally once the normal exists and once you start thinking about waits
and products.

Notation conventions across these distributions (Greek for parameters,
Roman with hats for estimators) will be respected throughout the book.
The hats become important in chapter 13, when we learn to estimate
parameters from data.

\end{tcolorbox}

\begin{tcolorbox}[enhanced jigsaw, arc=.35mm, bottomrule=.15mm, bottomtitle=1mm, breakable, colback=white, colbacktitle=quarto-callout-important-color!10!white, colframe=quarto-callout-important-color-frame, coltitle=black, left=2mm, leftrule=.75mm, opacityback=0, opacitybacktitle=0.6, rightrule=.15mm, title=\textcolor{quarto-callout-important-color}{\faExclamation}\hspace{0.5em}{Frontier note}, titlerule=0mm, toprule=.15mm, toptitle=1mm]

Modern probability is much wider than the six distributions above. Three
areas where the frontier is moving:

\emph{Heavy-tailed distributions and extreme value theory.} Many
real-world distributions (market returns, internet traffic, natural
disasters) have tails far heavier than the normal predicts. The field of
extreme value theory studies the maxima of random processes and provides
distributions (Gumbel, Fréchet, Weibull) that are the right tools when
the question is about rare extreme events rather than the bulk.

\emph{Distributions on networks and graphs.} The ``rectangle of rows and
columns'' assumption breaks down when the data is relational (who knows
whom, who trades with whom). Modern probability has been extended to
handle distributions on graphs, on permutations, on trees, on more
exotic combinatorial objects. The recurring trade panel of this book is
partly network data, and we will return to this in chapter 41.

\emph{Bayesian nonparametric models.} When the right distribution is
itself unknown, modern Bayesian methods place a prior over distributions
themselves (Dirichlet processes, Gaussian processes), letting the data
choose the right distributional family. This is one of the most active
areas in current statistics and machine learning.

For now, the six stories cover everything in Parts One and Two of the
book. The bigger zoo waits.

\end{tcolorbox}

\section{Practice problems}\label{practice-problems-2}

\begin{enumerate}
\def\labelenumi{\arabic{enumi}.}
\item
  For each variable below, name the distribution that most likely
  describes it: number of phone calls received in an hour; height of
  adult women in cm; whether a coin lands heads; income in a country;
  time between earthquakes; number of children per family.
\item
  A distribution has a long right tail. You take the log; the result
  looks normal. What family of distributions is the original?
\item
  A Poisson distribution has \(\lambda = 6.4\). What is the mean?
  Approximately what is the standard deviation? Why?
\item
  For Student's \(t\) at \(df = 5\), where do the tails sit relative to
  a standard normal? At \(df = 50\)?
\item
  Open the trade slice. Plot a histogram of
  \texttt{exports\_usd\_millions}. Then plot a histogram of
  \(\log(\text{exports_usd_millions})\). Which version is closer to a
  familiar shape from this chapter? What does that tell you about the
  underlying process?
\end{enumerate}

\begin{tcolorbox}[enhanced jigsaw, arc=.35mm, bottomrule=.15mm, bottomtitle=1mm, breakable, colback=white, colbacktitle=quarto-callout-caution-color!10!white, colframe=quarto-callout-caution-color-frame, coltitle=black, left=2mm, leftrule=.75mm, opacityback=0, opacitybacktitle=0.6, rightrule=.15mm, title=\textcolor{quarto-callout-caution-color}{\faFire}\hspace{0.5em}{Reflect before moving on}, titlerule=0mm, toprule=.15mm, toptitle=1mm]

\begin{itemize}
\tightlist
\item
  Why is ``what process produced this data'' a more useful question than
  ``what is the formula''?
\item
  Pick one distribution from the chapter. Explain its underlying story
  to a friend in three sentences.
\item
  The mean of a sample becomes normal as \(n\) grows, no matter what the
  parent distribution looks like. Why is this almost magical?
\end{itemize}

\end{tcolorbox}

\chapter{The grammar of graphics}\label{the-grammar-of-graphics}

\section{Anatomy of a chart}\label{anatomy-of-a-chart}

A chart is not a kind of thing. It is not ``a bar chart'' or ``a scatter
plot.'' A chart is a recipe with three layers, and once you see the
layers separately, the apparent variety of charts collapses into a much
smaller set of decisions.

The recipe was named by Leland Wilkinson in 1999 in \emph{The Grammar of
Graphics}, and refined by Hadley Wickham into the ggplot2 implementation
in R. Both versions agree on what the layers are.

\textbf{Data.} The dataset you are visualizing. Rectangular, of course.
Rows are observations, columns are variables.

\textbf{Aesthetics.} The mapping from columns of the dataset to visual
properties of the chart. Position on the x-axis. Position on the y-axis.
Color. Size. Shape. Each visual property is called an \emph{aesthetic},
and the chart-maker decides which dataset variable, if any, gets mapped
to each aesthetic.

\textbf{Geometry.} The kind of mark used to render each row. A point. A
bar. A line. A polygon.

That is the entire vocabulary. A scatter plot is \emph{data + position-x
+ position-y + point}. A bar chart is \emph{data + position-x + length-y
+ bar}. A line chart is \emph{data + position-x + position-y + line}. A
faceted small-multiples plot is the same with one extra aesthetic,
usually called \emph{facet}, that splits the data into panels.

\begin{figure*}

\centering{

\pandocbounded{\includegraphics[keepaspectratio]{prose/part1/../../pictures/figures/output/fig_grammar.pdf}}

}

\caption{\label{fig-grammar}Anatomy of a single chart, exploded into the
three layers of the grammar. The data is a small slice of trade flows.
One row maps to one geometric mark (a point). The aesthetics map columns
of the data to visual properties of the mark: x-position, y-position,
color. Change the geometry and you get a different chart from the same
data and aesthetics.}

\end{figure*}%

Once the grammar is in your head, ``what kind of chart should I make''
stops being a question about chart types and becomes a question about
which aesthetic each of your variables should get mapped to. That is the
right question. The answer depends on what you want the reader to see,
and on what the human visual system is good at.

\section{The perceptual hierarchy}\label{the-perceptual-hierarchy}

Not all aesthetics are created equal. The human visual system reads some
encodings much more accurately than others, and this fact has been
measured carefully.

The hierarchy, from most accurate to least, was established by Cleveland
and McGill in 1984 and confirmed in many studies since:

\begin{enumerate}
\def\labelenumi{\arabic{enumi}.}
\tightlist
\item
  Position along a common axis
\item
  Position along non-aligned axes
\item
  Length
\item
  Angle and slope
\item
  Area
\item
  Volume
\item
  Color hue and density
\end{enumerate}

This means that the same number, encoded as a position on a scatter
plot's y-axis, is read more accurately than the same number encoded as
the height of a bar (which is length), much more accurately than the
same number encoded as the area of a circle, and far more accurately
than the same number encoded as the lightness of a color.

\begin{figure*}

\centering{

\pandocbounded{\includegraphics[keepaspectratio]{prose/part1/../../pictures/figures/output/fig_perceptual_hierarchy.pdf}}

}

\caption{\label{fig-perceptual-hierarchy}The same five values rendered
five ways: as positions on a common axis, as bar lengths, as pie-slice
angles, as circle areas, and as color values. Cleveland and McGill's
experiments asked subjects to compare encoded values; accuracy fell off
in roughly that order. Position is read accurately. Color is read
poorly.}

\end{figure*}%

The implication is brutal and important. When you have a choice, encode
the variable you most want the reader to read accurately as a position.
When you do not have a choice (because position is already taken by
other variables), use length next. Use area only when ranking the values
matters but precise comparison does not. Use color hue last, and only
when the variable is categorical or the question is about rough
categories.

This is why bar charts are usually fine but pie charts are usually not
(the reader compares angles, item 4 on the list, when length on a bar
chart, item 3, would be more accurate). It is why bubble charts are
misleading (area, item 5). It is why heatmaps work for ``spot the hot
region'' but fail for ``compare exact values'' (color, the bottom of the
list).

\section{Small multiples}\label{small-multiples}

When you have one extra variable beyond what x and y are encoding, you
face a choice. Map the extra variable to a within-panel aesthetic
(color, shape, size), or use it to split the chart into panels.

Splitting into panels is called \emph{faceting}, or \emph{small
multiples}, and it has a strong claim to being the most underused
visualization technique in working analysis. The design move: instead of
overlaying everything on one chart, repeat the chart structure across
panels, with each panel showing a slice of the data.

\begin{figure*}

\centering{

\pandocbounded{\includegraphics[keepaspectratio]{prose/part1/../../pictures/figures/output/fig_faceting.pdf}}

}

\caption{\label{fig-faceting}A scatter of bilateral exports against
importer GDP, faceted by exporter country. Each panel uses the same axes
and the same encoding; only the panel-defining variable changes. The
pattern in any one panel is easy to read, and so is the contrast across
panels. The same data overlaid in one panel with color-by-country would
be a tangled mess.}

\end{figure*}%

Small multiples work because they use position (the panel grid is a
position encoding) for the splitting variable, and they keep the
within-panel encoding clean. They also support the eye in moving across
panels: if all panels share axes, the visual comparison is doing the
same kind of work as a scatter plot of summary statistics, except every
individual data point is visible.

When in doubt, facet.

\section{The POI palette as semantic
vocabulary}\label{the-poi-palette-as-semantic-vocabulary}

This book uses a fixed set of seven colors, and each one carries
meaning. The meanings are taught here and respected for the rest of the
book.

{\def\LTcaptype{none} % do not increment counter
\begin{longtable}[]{@{}
  >{\raggedright\arraybackslash}p{(\linewidth - 4\tabcolsep) * \real{0.3333}}
  >{\raggedright\arraybackslash}p{(\linewidth - 4\tabcolsep) * \real{0.3333}}
  >{\raggedright\arraybackslash}p{(\linewidth - 4\tabcolsep) * \real{0.3333}}@{}}
\toprule\noalign{}
\begin{minipage}[b]{\linewidth}\raggedright
Color
\end{minipage} & \begin{minipage}[b]{\linewidth}\raggedright
Hex
\end{minipage} & \begin{minipage}[b]{\linewidth}\raggedright
Meaning
\end{minipage} \\
\midrule\noalign{}
\endhead
\bottomrule\noalign{}
\endlastfoot
\textbf{INK} & \texttt{\#1a4f7a} & primary, default, the main story \\
\textbf{RUST} & \texttt{\#b85c38} & treatment, contrast, the alternative
perspective \\
\textbf{SAGE} & \texttt{\#5a7247} & control, comparison, the baseline \\
\textbf{GOLD} & \texttt{\#b8941e} & highlighted, derived, the
synthesized quantity \\
\textbf{VIOLET} & \texttt{\#6a5acd} & uncertain, predicted, the model's
guess \\
\textbf{DIM} & \texttt{\#8a8a8a} & background, de-emphasized, support \\
\textbf{TEAL} & \texttt{\#3a8a99} & alternative or secondary contrast \\
\end{longtable}
}

\begin{figure}

\centering{

\pandocbounded{\includegraphics[keepaspectratio]{prose/part1/../../pictures/figures/output/fig_palette.pdf}}

}

\caption{\label{fig-palette}The book's palette as swatches with semantic
labels. By chapter five you will recognize the colors without thinking.
A new figure with INK and RUST tells you, before you read the caption,
that there is a primary story and a contrast. A figure where INK fades
to DIM tells you which thing is being de-emphasized.}

\end{figure}%

This is unusual in textbooks. Most books pick colors per chapter, often
from a default palette that has no meaning attached. The result is that
the reader processes color afresh on every figure, which slows reading
and dilutes argument. The palette here is doing real work, and so are
you when you respect its conventions while making your own figures.

\section{Same data, four ways}\label{same-data-four-ways}

Here is one small dataset rendered as four different geometries. The
data is identical. The aesthetic mappings (which column maps to x and
which to y) are identical. Only the geometry changes.

\begin{figure*}

\centering{

\pandocbounded{\includegraphics[keepaspectratio]{prose/part1/../../pictures/figures/output/fig_same_data_four_ways.pdf}}

}

\caption{\label{fig-same-data}One small dataset rendered four ways:
scatter, bar, dot plot, slope graph. Each one is correct. Each one
privileges a different question. The scatter is good for ``how do x and
y relate.'' The bar chart is good for ``rank these from largest to
smallest.'' The dot plot is the same comparison without the visual
weight of bars. The slope graph is good for ``how did things change
between two points in time.''}

\end{figure*}%

Each one is correct. Each one is also, in a sense, wrong, because each
one privileges a particular kind of question. The lesson is not that one
of these is the best. The lesson is that the choice of geometry is a
substantive choice about which question the chart is designed to answer.

When you make a chart, ask first: what does the reader need to see? That
answer chooses the geometry. Everything else follows.

\begin{tcolorbox}[enhanced jigsaw, arc=.35mm, bottomrule=.15mm, bottomtitle=1mm, breakable, colback=white, colbacktitle=quarto-callout-warning-color!10!white, colframe=quarto-callout-warning-color-frame, coltitle=black, left=2mm, leftrule=.75mm, opacityback=0, opacitybacktitle=0.6, rightrule=.15mm, title=\textcolor{quarto-callout-warning-color}{\faExclamationTriangle}\hspace{0.5em}{Key insight}, titlerule=0mm, toprule=.15mm, toptitle=1mm]

A chart is data plus aesthetic mappings plus geometry, in that order.
Changing the geometry without changing the data answers a different
question. Changing the aesthetic mappings without changing the geometry
is what most ``chart improvement'' actually consists of. The perceptual
hierarchy (position beats length beats angle beats area beats color) is
the constraint that decides which mappings serve the reader and which do
not.

\end{tcolorbox}

\begin{tcolorbox}[enhanced jigsaw, arc=.35mm, bottomrule=.15mm, bottomtitle=1mm, breakable, colback=white, colbacktitle=quarto-callout-tip-color!10!white, colframe=quarto-callout-tip-color-frame, coltitle=black, left=2mm, leftrule=.75mm, opacityback=0, opacitybacktitle=0.6, rightrule=.15mm, title=\textcolor{quarto-callout-tip-color}{\faLightbulb}\hspace{0.5em}{Try this}, titlerule=0mm, toprule=.15mm, toptitle=1mm]

Pick any chart from a recent news article. Identify the three layers:
what is the dataset, which columns are mapped to which aesthetics, and
what is the geometry?

Then ask: would another geometry, with the same data and aesthetics,
answer the question better? Would a different aesthetic mapping (color
instead of x-position, say) work better, given what the reader is
supposed to see?

The exercise gets faster with practice. By the tenth chart you will be
doing it automatically.

\end{tcolorbox}

\begin{tcolorbox}[enhanced jigsaw, arc=.35mm, bottomrule=.15mm, bottomtitle=1mm, breakable, colback=white, colbacktitle=quarto-callout-note-color!10!white, colframe=quarto-callout-note-color-frame, coltitle=black, left=2mm, leftrule=.75mm, opacityback=0, opacitybacktitle=0.6, rightrule=.15mm, title=\textcolor{quarto-callout-note-color}{\faInfo}\hspace{0.5em}{Inside the math: a more formal grammar}, titlerule=0mm, toprule=.15mm, toptitle=1mm]

Wilkinson's \emph{Grammar of Graphics} (1999) decomposed a chart into
seven layers, of which the three above are the most consequential. The
full list adds \emph{scales} (how data values are mapped to visual
values), \emph{coordinates} (Cartesian, polar, geographic projections),
\emph{facets} (the small-multiples splitting variable), and
\emph{themes} (typography, gridlines, palette). Wickham's ggplot2
implementation in R follows the same structure.

The grammar is more general than this chapter's three-layer
simplification suggests. A polar coordinate system applied to a bar
chart produces a pie chart (with all the perceptual liabilities
described above). A logarithmic scale applied to a continuous axis
produces a log-log plot. A geographic projection applied to a position
aesthetic produces a map.

Each move is a transformation in the grammar, and the same chart is
described compactly by the sequence of transformations rather than by a
chart-type name. This is the conceptual move that makes ggplot2 readable
to people who have never seen it: the syntax mirrors the grammar.

\end{tcolorbox}

\begin{tcolorbox}[enhanced jigsaw, arc=.35mm, bottomrule=.15mm, bottomtitle=1mm, breakable, colback=white, colbacktitle=quarto-callout-important-color!10!white, colframe=quarto-callout-important-color-frame, coltitle=black, left=2mm, leftrule=.75mm, opacityback=0, opacitybacktitle=0.6, rightrule=.15mm, title=\textcolor{quarto-callout-important-color}{\faExclamation}\hspace{0.5em}{Frontier note}, titlerule=0mm, toprule=.15mm, toptitle=1mm]

The grammar of graphics is now several decades old and the field has
moved on. Three directions worth knowing.

\emph{Interactive grammars.} Vega and Vega-Lite extend the static
grammar to interaction (brushing, linked views, animation). The
visualizations on this book's website use Observable Plot, which is in
the same family. The interactive grammar is largely the static grammar
plus a vocabulary for events and selections.

\emph{Algebraic visualization grammars.} Recent work by Wills, Pu, and
others has formalized the grammar as an algebra, where chart
compositions can be reasoned about systematically. This is the basis for
tools that auto-recommend visualizations given the data.

\emph{Perceptual grammars.} Modern work continues to refine Cleveland
and McGill's hierarchy with new experimental data, especially for
unconventional encodings (animated transitions, 3D, AR/VR). Most of the
new findings reinforce the original hierarchy with caveats. Position
remains king.

The grammar is not finished, but the foundation is stable enough that
this book will rest on it without apology.

\end{tcolorbox}

\section{Practice problems}\label{practice-problems-3}

\begin{enumerate}
\def\labelenumi{\arabic{enumi}.}
\item
  Take any chart from a recent news article. Identify the data, the
  aesthetic mappings, and the geometry.
\item
  You have country, year, GDP, and continent. Sketch a chart that uses
  all four variables. Justify each encoding choice using the perceptual
  hierarchy.
\item
  Why is a pie chart usually a worse choice than a bar chart for the
  same comparison, even when both encode the same data?
\item
  Faceting moves an aesthetic from within-panel to position-of-panel.
  Explain why that is often an improvement, in terms of the perceptual
  hierarchy.
\item
  Take the trade slice. Make four versions of the same scatter
  (\texttt{exports} vs \texttt{distance}), encoding \texttt{exporter}
  four different ways: as color, as shape, as size, as facet. Which
  encoding lets you read the data fastest?
\end{enumerate}

\begin{tcolorbox}[enhanced jigsaw, arc=.35mm, bottomrule=.15mm, bottomtitle=1mm, breakable, colback=white, colbacktitle=quarto-callout-caution-color!10!white, colframe=quarto-callout-caution-color-frame, coltitle=black, left=2mm, leftrule=.75mm, opacityback=0, opacitybacktitle=0.6, rightrule=.15mm, title=\textcolor{quarto-callout-caution-color}{\faFire}\hspace{0.5em}{Reflect before moving on}, titlerule=0mm, toprule=.15mm, toptitle=1mm]

\begin{itemize}
\tightlist
\item
  What is one chart you have made or seen recently that violates the
  perceptual hierarchy? Could it be redrawn?
\item
  Color is at the bottom of the hierarchy but is used everywhere. Why?
  When is color actually justified?
\item
  Pick a question you have had to answer with data and name the right
  geometry for it.
\end{itemize}

\end{tcolorbox}

\chapter{Center, spread, shape}\label{center-spread-shape}

\section{Three numbers that summarize a
column}\label{three-numbers-that-summarize-a-column}

You have a column of numbers. A million rows, maybe. You want to
describe it in a way that fits in a sentence. What three numbers should
you reach for?

The answer in most introductory courses is: the mean, the standard
deviation, and not much else. That answer is not wrong. It is also not
enough. A column of numbers is described well by three things: where it
is centered, how spread out it is, and what shape it has. Each of these
can be measured several ways, and the choice of measurement is itself a
substantive decision.

The sections below lay out the choices. By the end of the chapter you
will be reaching for the right one without thinking, and the words
``mean'' and ``standard deviation'' will sound, properly, like just one
option among several.

\section{Center: mean, median, mode}\label{center-mean-median-mode}

The \emph{mean} is the arithmetic average. Sum the values, divide by the
count. It has a clean physical interpretation: if you put each data
value as a unit weight on a number line, the mean is the point where the
line balances.

The \emph{median} is the middle value when the data are sorted. Half the
values fall below it, half above. It has a different physical
interpretation: it is the point that minimizes the sum of absolute
distances to the data.

The \emph{mode} is the most common value, or for continuous data, the
value where the distribution peaks.

For perfectly symmetric distributions, all three are equal. For skewed
distributions, they pull apart. The mean chases the long tail. The
median ignores it. The mode plants itself at the peak.

\begin{figure}

\centering{

\pandocbounded{\includegraphics[keepaspectratio]{prose/part1/../../pictures/figures/output/fig_center_compared.pdf}}

}

\caption{\label{fig-center}A right-skewed distribution
(log-normal-shaped) with the three centers marked. The mode sits at the
peak, the median sits to its right (because more than half the area is
to the right of the mode), and the mean sits even further right (because
the long tail pulls the balance point that way). For symmetric
distributions, all three coincide.}

\end{figure}%

Income is the canonical example. A small number of very high earners
pull the mean far above the median. The median household income in a
country can be \$50,000 while the mean is \$80,000, simply because a
tiny fraction of the population earns enough to drag the average up.
When a journalist wants to make inequality look mild, they report the
mean. When they want to make it visible, they report the median. Both
numbers are correct. The choice of which to report is a choice about
what story to tell.

The general rule. For symmetric data, mean and median agree, and the
mean is more efficient (uses more of the data). For skewed data, the
median is the more honest summary of ``what does a typical observation
look like.'' The mean is still useful, but it tells you about the
\emph{aggregate} (the total divided by the count), not about the
\emph{typical}.

\section{Spread: range, variance, sd,
IQR}\label{spread-range-variance-sd-iqr}

A column with mean 50 could have all its values within a millimeter of
50, or it could have half its values at 0 and half at 100. The mean is
the same. The spread is wildly different.

Several ways to measure spread.

The \emph{range} is the largest value minus the smallest. Easy to
compute. Almost useless: a single outlier can blow the range up. Don't
use it for anything serious.

The \emph{variance} is the average of squared deviations from the mean,

\[
s^2 = \frac{1}{n-1} \sum_{i=1}^n (x_i - \bar{x})^2.
\]

The squaring matters. It prevents positive and negative deviations from
cancelling, and it weights large deviations more heavily than small
ones. Variance has the wrong units for direct interpretation (squared
whatever-the-data-was-in), so we usually report its square root, the
\emph{standard deviation} \(s\). Compare \(s\) to the mean to get a
sense of relative spread.

The \emph{interquartile range}, or IQR, is the spread between the 25th
percentile and the 75th percentile. Middle half of the data, as a single
number. By ignoring the top and bottom 25\%, the IQR becomes robust to
outliers in a way the standard deviation is not.

\begin{figure}

\centering{

\pandocbounded{\includegraphics[keepaspectratio]{prose/part1/../../pictures/figures/output/fig_spread_compared.pdf}}

}

\caption{\label{fig-spread}Three distributions with the same mean but
very different spreads. The left has standard deviation 0.5; the middle
has 1.0; the right has 2.5. The IQR ratios track the standard deviation
ratios because the data is normal-shaped. For non-normal data, the two
measures can diverge sharply.}

\end{figure}%

When the data is approximately normal, the standard deviation is the
right tool. The 68-95-99.7 rule from chapter 3 tells you exactly how to
interpret it. When the data is skewed or has outliers, the IQR is
usually more honest, because the standard deviation can be inflated by a
few extreme values that are not representative of the bulk.

\section{Shape: skewness and kurtosis without the
formulas}\label{shape-skewness-and-kurtosis-without-the-formulas}

Beyond center and spread, distributions have shape.

\emph{Skewness} measures asymmetry. A right-skewed distribution has a
long tail on the right (income, file sizes, earthquake magnitudes). A
left-skewed distribution has a long tail on the left (less common; one
example is age at death in low-mortality populations). A symmetric
distribution has zero skewness.

\emph{Kurtosis} measures tailedness, roughly speaking. A high-kurtosis
distribution has heavier tails than a normal: more probability of
extreme values. A low-kurtosis distribution has lighter tails. Financial
returns famously have high kurtosis: extreme moves are more common than
a normal distribution would predict, and the difference matters for risk
management.

\begin{figure*}

\centering{

\pandocbounded{\includegraphics[keepaspectratio]{prose/part1/../../pictures/figures/output/fig_shapes.pdf}}

}

\caption{\label{fig-shapes}Four shapes that recur: a roughly symmetric
(normal) distribution, a right-skewed (log-normal) distribution, a
heavy-tailed (Student's t with low df) distribution, and a bimodal
distribution. Knowing the formulas for skewness and kurtosis matters
less than recognizing the shapes.}

\end{figure*}%

The numerical formulas for skewness and kurtosis exist and you can look
them up. Knowing the formulas matters less than recognizing the shapes.
When you look at a histogram, your first move is: symmetric or skewed?
Light-tailed or heavy-tailed? One peak or several? Those four binary
questions catch most of the shape information you need for early-stage
analysis.

Bimodal distributions deserve special mention. When a distribution has
two peaks, the mean and median are both potentially meaningless: they
fall in the valley between the peaks, where almost no actual
observations live. Bimodality is usually a sign that two different
processes have been pooled together: men and women, treated and control,
before and after. The right move when you see bimodality is to find the
splitting variable and disaggregate.

\section{Robustness: what an outlier
does}\label{robustness-what-an-outlier-does}

An \emph{outlier} is an observation that lies far from the bulk. The
word does no semantic work; it is just a description of distance from
the center.

The mean and standard deviation are \emph{not robust} to outliers. A
single extreme value can change them substantially. The median and IQR
\emph{are robust}. A single extreme value barely moves them at all.

\begin{figure}

\centering{

\pandocbounded{\includegraphics[keepaspectratio]{prose/part1/../../pictures/figures/output/fig_robustness.pdf}}

}

\caption{\label{fig-robustness}Fifty points sampled near zero, with one
extreme outlier added at 100. The mean is dragged toward the outlier,
while the median barely moves. The standard deviation explodes; the IQR
holds steady. Robustness is the difference between summary statistics
that lie about an outlier and summary statistics that ignore it.}

\end{figure}%

The robustness distinction is more important than it sounds. Real
datasets have outliers. Sometimes those outliers are real (a billionaire
in an income survey). Sometimes they are errors (a typo, a unit
confusion, a sensor malfunction). In both cases, summary statistics that
are pulled around by the outliers will mislead. Robust statistics give
you a stable summary of the bulk, which is usually what you want before
you have figured out whether the outliers should be kept.

Use the median and IQR when you do not yet know what your outliers are.
Use the mean and standard deviation when the data is approximately
normal and you are confident the outliers are real and meaningful.

\section{The five-number summary and the
boxplot}\label{the-five-number-summary-and-the-boxplot}

A compact summary captures center, spread, and outliers in five numbers:
the minimum, the 25th percentile, the median, the 75th percentile, and
the maximum. Tukey called this the \emph{five-number summary}, and it is
one of the most useful descriptive tools in the discipline.

The \emph{boxplot} is the visualization of the five-number summary. A
box from the 25th to the 75th percentile, a line at the median, whiskers
extending to the most extreme non-outlier values, and individual points
for outliers beyond the whiskers (typically beyond 1.5 IQRs).

\begin{figure*}

\centering{

\pandocbounded{\includegraphics[keepaspectratio]{prose/part1/../../pictures/figures/output/fig_boxplot.pdf}}

}

\caption{\label{fig-boxplot}Anatomy of a boxplot, alongside boxplots for
the four shapes from earlier in the chapter. Boxplots are excellent for
comparing distributions across groups: row of boxplots, one per group,
is often the right first chart. They hide the within-box shape, which is
why a bimodal distribution and a uniform distribution can have similar
boxplots.}

\end{figure*}%

Boxplots are the right summary visualization for many comparisons across
groups: trade flows by year, test scores by class, response times by
condition. They show center and spread and outliers in a small
footprint, which makes them ideal for small multiples. A row of
boxplots, one per group, is often the right first chart for any
cross-group comparison.

Their weakness: they hide the \emph{shape} of the distribution within
the box. A bimodal distribution and a uniform distribution can have the
same boxplot. When shape matters, a \emph{violin plot} (which adds a
kernel density estimate to each side of the box) is the better choice.
We will meet violin plots formally in chapter seven.

\begin{tcolorbox}[enhanced jigsaw, arc=.35mm, bottomrule=.15mm, bottomtitle=1mm, breakable, colback=white, colbacktitle=quarto-callout-warning-color!10!white, colframe=quarto-callout-warning-color-frame, coltitle=black, left=2mm, leftrule=.75mm, opacityback=0, opacitybacktitle=0.6, rightrule=.15mm, title=\textcolor{quarto-callout-warning-color}{\faExclamationTriangle}\hspace{0.5em}{Key insight}, titlerule=0mm, toprule=.15mm, toptitle=1mm]

A column of numbers wants three things named: where it is centered, how
spread out it is, and what shape it has. Each can be measured several
ways, and the right choice depends on whether the data is symmetric,
whether outliers are present, and what story the summary needs to tell.
Robust statistics (median, IQR) are the default when you do not yet know
your data; classical statistics (mean, sd) come into their own when the
data is approximately normal.

\end{tcolorbox}

\begin{tcolorbox}[enhanced jigsaw, arc=.35mm, bottomrule=.15mm, bottomtitle=1mm, breakable, colback=white, colbacktitle=quarto-callout-tip-color!10!white, colframe=quarto-callout-tip-color-frame, coltitle=black, left=2mm, leftrule=.75mm, opacityback=0, opacitybacktitle=0.6, rightrule=.15mm, title=\textcolor{quarto-callout-tip-color}{\faLightbulb}\hspace{0.5em}{Try this}, titlerule=0mm, toprule=.15mm, toptitle=1mm]

Open the bilateral trade panel. For the \texttt{exports\_usd\_millions}
column, compute: mean, median, standard deviation, IQR, and the
five-number summary. Plot a histogram and a boxplot.

Which summary best describes the column? If you had to give one number
to a journalist who wanted to know ``how much typical bilateral trade is
between these countries,'' which would you give?

Then take the log of \texttt{exports\_usd\_millions} and recompute.
Notice which summary changes most dramatically and which barely changes.
The answer tells you something about the underlying distribution and
about robustness, both at once.

\end{tcolorbox}

\begin{tcolorbox}[enhanced jigsaw, arc=.35mm, bottomrule=.15mm, bottomtitle=1mm, breakable, colback=white, colbacktitle=quarto-callout-note-color!10!white, colframe=quarto-callout-note-color-frame, coltitle=black, left=2mm, leftrule=.75mm, opacityback=0, opacitybacktitle=0.6, rightrule=.15mm, title=\textcolor{quarto-callout-note-color}{\faInfo}\hspace{0.5em}{Inside the math: where the n-1 comes from}, titlerule=0mm, toprule=.15mm, toptitle=1mm]

The variance formula has \(n-1\) in the denominator, not \(n\). This is
\emph{Bessel's correction}, and the reason for it is one of the cleanest
derivations in statistics.

When you estimate the population variance from a sample, you use the
sample mean \(\bar{x}\) in place of the unknown population mean \(\mu\).
The sample mean minimizes \(\sum(x_i - c)^2\) over choices of \(c\),
which means \(\sum(x_i - \bar{x})^2\) is \emph{systematically smaller}
than \(\sum(x_i - \mu)^2\) would have been. Dividing by \(n\) would
therefore underestimate the population variance.

The correction is exactly \(n/(n-1)\): dividing by \(n-1\) instead of
\(n\) produces an unbiased estimator of the population variance. The
intuition is that we have used up ``one degree of freedom'' in
estimating the mean from the same data, so the effective sample size for
estimating spread is \(n-1\) rather than \(n\).

This shows up everywhere in inferential statistics. Whenever you see
\(n-1\) in a denominator, look for an estimated mean somewhere in the
numerator. Degrees of freedom are the bookkeeping system for this
correction, and we will encounter them again in the t distribution
(chapter 13), in regression (chapter 16), and in everything that
follows.

\end{tcolorbox}

\begin{tcolorbox}[enhanced jigsaw, arc=.35mm, bottomrule=.15mm, bottomtitle=1mm, breakable, colback=white, colbacktitle=quarto-callout-important-color!10!white, colframe=quarto-callout-important-color-frame, coltitle=black, left=2mm, leftrule=.75mm, opacityback=0, opacitybacktitle=0.6, rightrule=.15mm, title=\textcolor{quarto-callout-important-color}{\faExclamation}\hspace{0.5em}{Frontier note}, titlerule=0mm, toprule=.15mm, toptitle=1mm]

Robust statistics is a substantial sub-field of modern statistics.
Beyond the median and IQR, there are robust analogues of almost every
classical procedure: M-estimators, trimmed means, winsorized means, the
median absolute deviation (MAD), Huber regression, and many more. The
frontier work is in \emph{high-dimensional} robustness: when the data
has thousands of variables, classical procedures break down in ways that
simple median-and-IQR replacements do not fix.

In modern empirical economics, the practice of reporting both classical
and robust summaries (or robust standard errors, which we will meet in
chapter 34) is now standard for any serious paper. The reason is partly
the credibility revolution and partly the recognition that real economic
data has fat tails, outliers, and clusters, and pretending otherwise
produces fragile inference.

\end{tcolorbox}

\section{Practice problems}\label{practice-problems-4}

\begin{enumerate}
\def\labelenumi{\arabic{enumi}.}
\item
  Compute by hand: mean, median, mode, range, variance, sd, IQR, and the
  five-number summary for the values
  \(\{1, 3, 3, 5, 7, 8, 9, 12, 100\}\).
\item
  Which summary statistic from the previous problem changes the most if
  you remove the 100? Which barely changes? What does that tell you
  about robustness?
\item
  A distribution has skewness \(+2\). What does that tell you about the
  shape? What if the skewness is \(-2\)?
\item
  A boxplot shows three distributions side by side. Without seeing
  histograms, which features can you compare confidently? Which would
  you want histograms for?
\item
  Open the trade slice. Compute the mean and the median of
  \texttt{exports\_usd\_millions}. Which is larger? Why? Which would you
  report to a journalist?
\end{enumerate}

\begin{tcolorbox}[enhanced jigsaw, arc=.35mm, bottomrule=.15mm, bottomtitle=1mm, breakable, colback=white, colbacktitle=quarto-callout-caution-color!10!white, colframe=quarto-callout-caution-color-frame, coltitle=black, left=2mm, leftrule=.75mm, opacityback=0, opacitybacktitle=0.6, rightrule=.15mm, title=\textcolor{quarto-callout-caution-color}{\faFire}\hspace{0.5em}{Reflect before moving on}, titlerule=0mm, toprule=.15mm, toptitle=1mm]

\begin{itemize}
\tightlist
\item
  Mean and median tell you different things. Pick a real-world quantity
  where you would report the median over the mean. Explain.
\item
  The IQR ignores the top and bottom 25\%. What information are you
  giving up by reporting it instead of the standard deviation?
\item
  The bimodality warning sign comes up again and again. Why is
  bimodality usually a clue to disaggregate?
\end{itemize}

\end{tcolorbox}

\chapter{Lying with charts}\label{lying-with-charts}

\section{A defensive reading skill}\label{a-defensive-reading-skill}

This is a different kind of chapter. The previous chapters built up the
visual vocabulary the book uses. This one teaches you to read other
people's charts critically, because most charts you encounter outside of
textbooks are made by people with an argument they want to win, and a
well-made chart is a powerful tool for winning arguments dishonestly.

The techniques below are not separate problems. They are the same
problem (a mismatch between visual encoding and the truth of the
underlying data) applied to different aesthetics. Truncated y-axes lie
about position. Cherry-picked windows lie about what is in scope. Area
encodings lie when length is what is being compared. Dual axes lie about
what corresponds to what. Log scales lie about effect sizes. Each one
will fool you the first time you see it. Once you can name the trick,
you stop being fooled.

The chapter ends with a short list of honest defaults that prevent most
of these errors when you make charts of your own. Read them last.

\section{The truncated y-axis}\label{the-truncated-y-axis}

The most common visual lie. A bar chart shows that quantity X went from
100 to 102, but the y-axis starts at 99 and ends at 103. Visually, the
second bar looks much taller than the first: maybe four times taller.
Numerically, X grew by 2\%.

\begin{figure*}

\centering{

\pandocbounded{\includegraphics[keepaspectratio]{prose/part1/../../pictures/figures/output/fig_truncated_axis.pdf}}

}

\caption{\label{fig-truncated}Same data, two y-axes. On the left, the
y-axis starts at zero, and a 2\% growth looks like a 2\% growth. On the
right, the y-axis starts at 99, and the same 2\% growth looks like the
second value is several times larger than the first. The data is
identical. The visual story is opposite.}

\end{figure*}%

The defensive move: always check where the y-axis starts. If it does not
start at zero on a bar chart of magnitudes, the chart is amplifying or
muting differences relative to a chart that does start at zero.
Sometimes that amplification is editorially fine (small changes in CPI
matter; you would not want to start at zero for a chart of historical
CPI, because zero CPI is a meaningless reference). More often it is
dishonest.

A useful rule. Bar charts should almost always start at zero, because
the visual story of a bar chart is ``compare the heights.'' Line charts
are different. Line charts about percentage changes or rates can
legitimately start anywhere, because the visual story is about
\emph{changes} over time, not absolute heights. Be especially suspicious
of bars that do not start at zero.

\section{Cherry-picked time windows}\label{cherry-picked-time-windows}

The same time series can tell completely different stories depending on
which part of it you show. Pick a starting point near a peak and you get
a downward trend. Pick a starting point near a trough and you get an
upward one. Both are honest in the sense that the data shown is real.
Both are dishonest in the sense that the choice of window is the
editorial argument, and a chart that does not show the longer context is
hiding the argument behind a curated frame.

\begin{figure*}

\centering{

\pandocbounded{\includegraphics[keepaspectratio]{prose/part1/../../pictures/figures/output/fig_cherry_pick.pdf}}

}

\caption{\label{fig-cherry-pick}A 30-year time series with a clear
upward trend. On the left, the full series. On the right, a five-year
window selected from the same series, chosen to suggest decline. Both
panels show real data. Only one tells the truth about what is going on.}

\end{figure*}%

The defensive move: ask why this window. If the chart shows five years
and you suspect more data exists, look for it. If the chart starts at a
suspiciously specific date, check what was happening just before.
Markets, climate series, public-health metrics, all of these have
natural reference points (a peak, a recession, a prior policy change)
that are deliberately included or excluded to shape the story.

\section{Area when length was needed}\label{area-when-length-was-needed}

A bar twice as tall has twice the height. A circle twice the diameter
has four times the area. A sphere twice the diameter has eight times the
volume. When a chart encodes a quantity as the area of a circle (a
``bubble chart'' or ``infographic'') or as the volume of a 3D shape, the
visual sense of magnitude grows faster than the underlying number, and
the reader systematically overestimates differences.

\begin{figure}

\centering{

\pandocbounded{\includegraphics[keepaspectratio]{prose/part1/../../pictures/figures/output/fig_area_distortion.pdf}}

}

\caption{\label{fig-area-distortion}Two bars representing 100 and 200.
On the left, only the height changes (the honest encoding). On the
right, the second bar is also drawn twice as wide, so its area is four
times larger than the first bar's, even though the underlying number is
only twice as large. The reader's eye overestimates by a factor of two.}

\end{figure}%

This goes wrong in two ways in the wild. One is genuine area encoding
(``this circle represents 100 million people; this circle, twice as
large in area, represents 200 million''). The reader's eye does not
weigh area linearly, so the ``twice as large'' judgment is unreliable.
The other is \emph{unintended} area encoding: a bar chart where the bars
get wider as well as taller, or a 3D bar chart where perspective makes
the back bars look different from the front ones at the same height.
Both are common in business presentations.

The defensive move: when you see area or volume as an encoding, do the
calculation. The 2x circle has 4x the area; ask whether the underlying
number is 2x or 4x.

\section{Dual-axis chicanery}\label{dual-axis-chicanery}

Two time series on the same chart, with two y-axes, one on each side.
Each axis is scaled differently. The two series appear to track each
other beautifully, suggesting a relationship.

The trick is in the axis scaling. The author has chosen the two axis
ranges to make the curves overlap. Different ranges would produce no
apparent relationship. The eye, however, reads ``they are tracking each
other'' with no awareness that the relationship is an artifact of axis
choice.

\begin{figure}

\centering{

\pandocbounded{\includegraphics[keepaspectratio]{prose/part1/../../pictures/figures/output/fig_dual_axis.pdf}}

}

\caption{\label{fig-dual-axis}A dual-axis chart of two unrelated time
series. The left axis runs from 0 to 100; the right axis runs from 5 to
7. The two curves appear to track each other across the visible range,
suggesting a strong correlation. They have none. Slide the right axis up
or down and the apparent relationship disappears.}

\end{figure}%

Dual-axis charts are not always dishonest. Sometimes the author
genuinely wants to show two related quantities measured in different
units (temperature and humidity, say). But the burden of proof should
fall on the author. If the two series are conceptually unrelated, dual
axes will manufacture an apparent correlation that does not exist.

The defensive move: separate the two series into two stacked panels with
their own axes. The relationship, if real, will still be visible. The
relationship, if manufactured by axis manipulation, will dissolve.

\section{Log scales that hide effect
sizes}\label{log-scales-that-hide-effect-sizes}

The previous tricks amplify small effects. Log scales mute large ones. A
series that grows from 1 to 1,000 takes a giant vertical step on a
linear scale and a small one on a log scale. On the log scale, the
growth from 1 to 10 looks the same as the growth from 100 to 1,000, even
though both represent tenfold increases. The constant log-step is good
for showing rates of growth, but it can be used to make explosive growth
look gentle.

\begin{figure*}

\centering{

\pandocbounded{\includegraphics[keepaspectratio]{prose/part1/../../pictures/figures/output/fig_log_hides.pdf}}

}

\caption{\label{fig-log-hides}An exponentially growing series shown on a
linear axis (left) and a logarithmic axis (right). The linear chart
screams. The log chart shrugs. Both show the same data. The choice of
scale is the rhetorical move.}

\end{figure*}%

A common honest use: financial returns over decades, where the
eye-relevant quantity is the percentage change, not the dollar level. A
common dishonest use: showing a runaway quantity on a log scale and
pointing to its ``modest growth,'' when the underlying number has gone
up by several orders of magnitude.

The defensive move: when you see a log axis, mentally re-render the data
on a linear axis. If the linear version would show an alarming surge,
the log scale was muting an honest signal.

\section{The honest defaults}\label{the-honest-defaults}

Five rules that prevent most of the lies above when you are the one
making the chart.

Bars start at zero. Line charts can start anywhere if the reader is
told. Time-series windows include the most relevant context, not just
the part that supports the argument. Quantities are encoded as length or
position, not as area or volume, unless the question is genuinely about
area or volume. Two unrelated series live in two stacked panels, not on
dual axes. Log scales are used when growth rate is the question, and
never when the level is.

That is the entire defensive checklist. Apply it to every chart in the
news this week. You will be surprised how many fail.

\begin{tcolorbox}[enhanced jigsaw, arc=.35mm, bottomrule=.15mm, bottomtitle=1mm, breakable, colback=white, colbacktitle=quarto-callout-warning-color!10!white, colframe=quarto-callout-warning-color-frame, coltitle=black, left=2mm, leftrule=.75mm, opacityback=0, opacitybacktitle=0.6, rightrule=.15mm, title=\textcolor{quarto-callout-warning-color}{\faExclamationTriangle}\hspace{0.5em}{Key insight}, titlerule=0mm, toprule=.15mm, toptitle=1mm]

Every visual lie in this chapter is the same kind of lie: a mismatch
between the visual encoding and the underlying number. Position lies
(truncated axes), scope lies (cherry-picked windows), area lies (bubble
charts), correspondence lies (dual axes), and rate-vs-level lies (log
scales). Five tricks. Once named, they stop fooling you.

\end{tcolorbox}

\begin{tcolorbox}[enhanced jigsaw, arc=.35mm, bottomrule=.15mm, bottomtitle=1mm, breakable, colback=white, colbacktitle=quarto-callout-tip-color!10!white, colframe=quarto-callout-tip-color-frame, coltitle=black, left=2mm, leftrule=.75mm, opacityback=0, opacitybacktitle=0.6, rightrule=.15mm, title=\textcolor{quarto-callout-tip-color}{\faLightbulb}\hspace{0.5em}{Try this}, titlerule=0mm, toprule=.15mm, toptitle=1mm]

Find a chart in a recent news article that you have an opinion about
(politically, economically, scientifically, whatever). Apply the
five-trick checklist. Does the chart truncate the y-axis? Cherry-pick a
window? Use area when length would have done? Stack two unrelated series
on dual axes? Use a log scale to hide a level?

Now do the harder version: find a chart in a recent paper from a journal
you respect. Apply the same checklist. Find at least one thing the chart
could be doing better. Many published charts do not pass the checklist;
the journal-name does not absolve them.

The exercise is not about catching dishonesty (most authors are not
malicious). It is about training your eye to notice when a visual is
doing rhetorical work that the underlying data does not support.

\end{tcolorbox}

\begin{tcolorbox}[enhanced jigsaw, arc=.35mm, bottomrule=.15mm, bottomtitle=1mm, breakable, colback=white, colbacktitle=quarto-callout-note-color!10!white, colframe=quarto-callout-note-color-frame, coltitle=black, left=2mm, leftrule=.75mm, opacityback=0, opacitybacktitle=0.6, rightrule=.15mm, title=\textcolor{quarto-callout-note-color}{\faInfo}\hspace{0.5em}{Inside the math: why the eye fails on area}, titlerule=0mm, toprule=.15mm, toptitle=1mm]

Stevens' law, from psychophysics, says that the perceived intensity of a
stimulus follows a power law: \(S = k \cdot I^a\), where \(I\) is the
physical intensity and \(a\) is an exponent that depends on the stimulus
type. For length, \(a\) is approximately 1.0 (the eye reads length
almost linearly). For area, \(a\) is approximately 0.7 (the eye
underestimates large areas relative to small ones). For volume, \(a\) is
around 0.6.

The implication: when you encode a quantity as area, the reader
systematically misjudges large values. A circle that is 4x larger in
area is judged to be only \(4^{0.7} \approx 2.6\) times larger by eye.
The number is right; the perception is wrong.

This is the deep reason area-based charts fail at precise comparison,
and it is why position and length (where the eye reads linearly) sit at
the top of the perceptual hierarchy. The hierarchy is not a stylistic
preference; it is a description of human perception, derived from
controlled experiments going back to Cleveland and McGill in the 1980s
and Stevens in the 1950s.

\end{tcolorbox}

\begin{tcolorbox}[enhanced jigsaw, arc=.35mm, bottomrule=.15mm, bottomtitle=1mm, breakable, colback=white, colbacktitle=quarto-callout-important-color!10!white, colframe=quarto-callout-important-color-frame, coltitle=black, left=2mm, leftrule=.75mm, opacityback=0, opacitybacktitle=0.6, rightrule=.15mm, title=\textcolor{quarto-callout-important-color}{\faExclamation}\hspace{0.5em}{Frontier note}, titlerule=0mm, toprule=.15mm, toptitle=1mm]

A growing literature studies \emph{visualization deception} as a
research topic. Pandey and others (2015) showed experimentally that
truncated axes, area encodings, and cherry-picked windows shift readers'
beliefs about underlying data, even when the readers know the chart
might be biased. The effect persists for trained scientists.

The policy response, in some venues, is the requirement of \emph{open
data and open code} alongside any chart in a published paper. If readers
can re-render the data themselves with honest defaults, the editorial
freedom to mislead is sharply reduced. Journals like \emph{AEA: Applied
Economics} and the broader open-science movement have moved in this
direction. The book's worldview, if you have not noticed, is sympathetic
to it.

For a longer treatment of this material, see Tufte's \emph{The Visual
Display of Quantitative Information} (1983), which remains the standard
reference. Tufte's ``lie factor'' formalizes the area distortion in this
chapter into a single number you can compute for any chart.

\end{tcolorbox}

\section{Practice problems}\label{practice-problems-5}

\begin{enumerate}
\def\labelenumi{\arabic{enumi}.}
\item
  Find one chart in a news article from this week. Apply the five-trick
  checklist. What does the chart do well? What does it hide?
\item
  Tufte defined the \emph{lie factor} as the ratio of the visual size of
  a change to the actual size. Compute it for a chart you find in the
  wild.
\item
  When is it legitimate to truncate a y-axis? Give two examples and
  explain what makes them legitimate.
\item
  Two unrelated time series look correlated on a dual-axis chart. Sketch
  the same data on stacked panels with separate axes. What changes?
\item
  A log scale is honest when growth rate is the question and dishonest
  when level is. Find one example of each in current news and defend
  your call.
\end{enumerate}

\begin{tcolorbox}[enhanced jigsaw, arc=.35mm, bottomrule=.15mm, bottomtitle=1mm, breakable, colback=white, colbacktitle=quarto-callout-caution-color!10!white, colframe=quarto-callout-caution-color-frame, coltitle=black, left=2mm, leftrule=.75mm, opacityback=0, opacitybacktitle=0.6, rightrule=.15mm, title=\textcolor{quarto-callout-caution-color}{\faFire}\hspace{0.5em}{Reflect before moving on}, titlerule=0mm, toprule=.15mm, toptitle=1mm]

\begin{itemize}
\tightlist
\item
  Which of the five tricks have you accidentally used in a chart of your
  own?
\item
  The chapter argues that most visual lies are mismatches between
  encoding and data. Can you think of a sixth kind of mismatch that was
  not named?
\item
  What is the strongest argument \emph{for} a misleading chart? When (if
  ever) is editorial framing legitimate?
\end{itemize}

\end{tcolorbox}

\chapter{Underused visualizations}\label{underused-visualizations}

\section{Beyond the standard
repertoire}\label{beyond-the-standard-repertoire}

The standard chart repertoire in introductory courses is small: bars,
lines, scatters, pies, histograms. Boxplots, if the instructor is
feeling generous. Most students leave their first stats course with that
as their entire vocabulary, and most working analysts add at most one or
two more over the course of a career.

This is a poverty. The visualization literature has produced many chart
types that answer specific questions better than anything in the
standard repertoire. Knowing them is the difference between an analyst
who reaches for a bar chart in every situation and one who reaches for
the right tool. The chart types in this chapter are not exotic. They are
well-understood, available in every plotting library, and
disproportionately useful for the kinds of questions empirical research
actually asks. They are simply not taught.

Six chart types follow. Each gets a section and a question. Read the
figures first, the prose second.

\section{The slope graph}\label{the-slope-graph}

Tufte's slope graph is the answer to ``how did things change between two
points in time?'' It is two columns of numbers connected by lines: the
value of each unit at time A on the left, the value at time B on the
right, a line connecting them. Lines that go up represent increases.
Lines that go down represent decreases. The eye reads the slope of each
line to compare changes across units.

\begin{figure}

\centering{

\pandocbounded{\includegraphics[keepaspectratio]{prose/part1/../../pictures/figures/output/fig_slope_graph.pdf}}

}

\caption{\label{fig-slope}Bilateral export volumes from each of our six
exporters to the United States, in 2014 and 2023, connected by lines.
Lines that rise represent growth; lines that fall represent decline. The
slope graph encodes both the level at each endpoint and the direction of
change in a single visual unit, which a bar chart of differences cannot
do.}

\end{figure}%

Why this works better than alternatives. A bar chart of differences
answers the same question, but the slope graph also shows the levels at
the start and end, which provides context. Two bar charts side by side
answer the same question, but the eye has to do work to track which unit
went where. The slope graph encodes change as the slope of a line; what
the eye is asked to read is differences in slope, which is a much easier
comparison than absolute angle.

When to reach for it. Two-point comparisons across many units. Survey
results before and after an intervention. Pollutant levels at the
beginning and end of a regulation period. Country-level outcomes across
two policy regimes. The slope graph is at its best with somewhere
between five and thirty units. Fewer is overkill; more becomes
spaghetti.

\section{The Cleveland dot plot}\label{the-cleveland-dot-plot}

A Cleveland dot plot is a bar chart with the bars replaced by single
dots. Each dot sits at the value, with a thin guideline back to a
reference (often zero, sometimes the median). That is the entire chart.

\begin{figure}

\centering{

\pandocbounded{\includegraphics[keepaspectratio]{prose/part1/../../pictures/figures/output/fig_dot_plot.pdf}}

}

\caption{\label{fig-dot}Total exports from each of our six countries in
2023, ranked. Each dot is the value; the thin guideline traces back to
zero. The same data as a bar chart would use far more ink to encode the
same information, with no gain in accuracy of reading.}

\end{figure}%

Why this is better than alternatives. Bar charts use ink to encode the
same value twice: once as the position of the bar's top, and once as the
bar's length. The eye reads position more accurately than length (per
the perceptual hierarchy in chapter 4), so the body of the bar is
decoration, not information. Dot plots strip the decoration. The result
is more readable, especially when the number of categories is large.

The Cleveland dot plot is also better when the data does not start at
zero, where bar charts become misleading (chapter 6). Dots can sit
anywhere on the axis without implying a magnitude from the visual length
of a bar. Stock prices, temperature anomalies, log-transformed values:
all comfortable in dot plots, awkward in bars.

When to reach for it. Ranked comparison across many categories,
especially when zero is not a meaningful reference. Country rankings,
journal impact factors, expression levels across genes, performance
benchmarks across teams.

\section{The hex bin}\label{the-hex-bin}

A scatter plot stops working when there are too many points. With a
thousand points, you can see structure. With a million, the points pile
on top of each other and you see a black blob. The hex bin solves this.

A hex bin chart divides the plotting area into hexagonal cells and
counts the number of points falling into each cell. Each cell is then
colored by its count. Dense regions appear dark; sparse regions appear
pale. The result is a heatmap of point density that reveals structure
invisible in the underlying scatter.

\begin{figure}

\centering{

\pandocbounded{\includegraphics[keepaspectratio]{prose/part1/../../pictures/figures/output/fig_hex_bin.pdf}}

}

\caption{\label{fig-hex}A hex-binned scatter of 100,000 points drawn
from a slightly correlated noisy joint distribution. A regular scatter
at this point count is a black smudge; the hex bin reveals the
underlying density gradient and the modes. Hexagons (rather than
squares) tile the plane more isotropically, so visual artifacts from the
grid itself are minimized.}

\end{figure}%

Why this works better than alternatives. With many points, a regular
scatter is unreadable. A 2D histogram with rectangular bins is a step up
but suffers from the rectangular grid showing through visually. Hexagons
are nearly circular, which means they bin similarly in all directions,
and they tile the plane without leaving gaps. The aesthetic result is
smoother than rectangular bins and more accurate than a smoothed density
plot for finding precise modes.

When to reach for it. Any scatter plot with more than about 5,000
points. Astronomy data, genomic data, large survey or transactional
data, log files. If the underlying scatter would be a black blob, the
hex bin is your friend.

\section{The ridge plot}\label{the-ridge-plot}

A ridge plot, sometimes called a joy plot, is a stacked sequence of
small density plots, one per category, arranged vertically with each
density slightly overlapping the one below. The result looks a bit like
a topographic map of distributions.

\begin{figure}

\centering{

\pandocbounded{\includegraphics[keepaspectratio]{prose/part1/../../pictures/figures/output/fig_ridge_plot.pdf}}

}

\caption{\label{fig-ridge}Distributions of bilateral export volumes (log
scale) across the years 2014 to 2023 in our trade slice. One density per
year, stacked from oldest at the top to most recent at the bottom. The
visual story is whether and how the distribution drifts over time, which
would be impossible to read from a single overlaid histogram and tedious
to read from a row of small multiples.}

\end{figure}%

Why this works better than alternatives. To compare the distribution of
a quantity across many categories (years, countries, conditions), the
alternatives are: many small histograms (small multiples), violin plots
side by side, or a single histogram colored by category (which becomes
mush). The ridge plot is the most space-efficient of these for showing
many distributions at once. The eye reads the shapes top to bottom and
immediately spots which categories have unusual distributions.

When to reach for it. Comparing the shape of a distribution across five
to fifty categories. Time-evolution of a distribution. Distribution of
one variable across many groups. The ridge plot is what you want when
the comparison is shape-vs-shape, not
summary-statistic-vs-summary-statistic.

\section{The mosaic plot}\label{the-mosaic-plot}

A mosaic plot is the visualization of a contingency table. The plot is a
rectangle divided proportionally by one variable along one axis. Each
resulting strip is then divided by another variable along the other
axis. The areas of the resulting tiles are proportional to the joint
counts.

\begin{figure}

\centering{

\pandocbounded{\includegraphics[keepaspectratio]{prose/part1/../../pictures/figures/output/fig_mosaic.pdf}}

}

\caption{\label{fig-mosaic}A mosaic plot of 2,000 simulated passenger
records by class (1st, 2nd, 3rd) and survival outcome. Each rectangle's
area is proportional to the count in that cell. The widths of the
columns encode the marginal distribution of class. The heights inside
each column encode the conditional distribution of survival within that
class. Departures from independence (different conditional heights
across columns) are visible as misalignments of the horizontal cuts.}

\end{figure}%

Why this works better than alternatives. A grouped bar chart of a
contingency table answers questions about counts, but it does not show
the marginal distribution of each variable, and it makes the joint
distribution hard to read because the bars within a group are at
different heights. The mosaic plot encodes everything at once: the
marginal distribution along one axis, the conditional distribution
within each level along the other, and the joint counts as areas. The
eye spots which combinations are over- or under-represented relative to
independence.

When to reach for it. Two or three categorical variables and you want to
see association. Survival rates by class on a ship. Voting patterns by
demographic and region. Confusion matrices in classification when you
want to see joint structure beyond the diagonal.

\section{The alluvial diagram}\label{the-alluvial-diagram}

An alluvial diagram, in its directed-flow form often called a Sankey,
shows how a population redistributes between categorical states. Stripes
leave one set of categories on the left and arrive at another set on the
right, with the width of each stripe proportional to the number of units
making that transition.

\begin{figure}

\centering{

\pandocbounded{\includegraphics[keepaspectratio]{prose/part1/../../pictures/figures/output/fig_alluvial.pdf}}

}

\caption{\label{fig-alluvial}Trade-partner rank changes between 2014 and
2023 for our six countries. Each country occupies a stripe whose
vertical position encodes its rank as an importer of bilateral exports.
Stripes flow from left (2014 rank) to right (2023 rank), with width
proportional to trade volume in 2023. Crossings reveal countries that
climbed or fell in the rankings.}

\end{figure}%

Why this works better than alternatives. Tracking flow between
categories is hard with bars or lines. A grouped bar chart of ``before''
and ``after'' loses the connection between them. A Sankey or alluvial
diagram makes the connection explicit: the visual width of each stripe
is the number of units following that path.

When to reach for it. Categorical transitions over time or between
conditions. Voter movement between elections. Customer-state movement
(free → paid → churned). Trade partner reassignments. Employment-status
changes. The alluvial is the right tool whenever the question is ``where
did they go?''

\begin{tcolorbox}[enhanced jigsaw, arc=.35mm, bottomrule=.15mm, bottomtitle=1mm, breakable, colback=white, colbacktitle=quarto-callout-warning-color!10!white, colframe=quarto-callout-warning-color-frame, coltitle=black, left=2mm, leftrule=.75mm, opacityback=0, opacitybacktitle=0.6, rightrule=.15mm, title=\textcolor{quarto-callout-warning-color}{\faExclamationTriangle}\hspace{0.5em}{Key insight}, titlerule=0mm, toprule=.15mm, toptitle=1mm]

The standard chart repertoire teaches answers to a small set of
questions. Each chart type in this chapter is the answer to a question
the standard repertoire handles poorly. Knowing the right tool for the
question is the difference between an analyst who fights their data and
one who reads it.

\end{tcolorbox}

\begin{tcolorbox}[enhanced jigsaw, arc=.35mm, bottomrule=.15mm, bottomtitle=1mm, breakable, colback=white, colbacktitle=quarto-callout-tip-color!10!white, colframe=quarto-callout-tip-color-frame, coltitle=black, left=2mm, leftrule=.75mm, opacityback=0, opacitybacktitle=0.6, rightrule=.15mm, title=\textcolor{quarto-callout-tip-color}{\faLightbulb}\hspace{0.5em}{Try this}, titlerule=0mm, toprule=.15mm, toptitle=1mm]

For each scenario below, name the chart type from this chapter that fits
best and explain why in one sentence:

\begin{enumerate}
\def\labelenumi{\arabic{enumi}.}
\tightlist
\item
  You have GDP per capita for 30 countries in 1990 and 2020. You want to
  show which countries gained ground and which lost it.
\item
  You have a million bilateral trade flows and want to see the joint
  shape of GDP and distance.
\item
  You have a survey with five Likert items broken down by twelve
  demographic groups. You want to show how response distributions differ
  across groups.
\item
  You have students moving between four academic majors over four years
  of college. You want to show the dominant trajectories.
\item
  You have voting patterns by income decile and education level. You
  want to show which combinations vote which way.
\end{enumerate}

After answering, draw at least one of these in a tool of your choice.
The hardest part of using non-standard charts is reaching for them; once
you do, the implementation is rarely the bottleneck.

\end{tcolorbox}

\begin{tcolorbox}[enhanced jigsaw, arc=.35mm, bottomrule=.15mm, bottomtitle=1mm, breakable, colback=white, colbacktitle=quarto-callout-note-color!10!white, colframe=quarto-callout-note-color-frame, coltitle=black, left=2mm, leftrule=.75mm, opacityback=0, opacitybacktitle=0.6, rightrule=.15mm, title=\textcolor{quarto-callout-note-color}{\faInfo}\hspace{0.5em}{Inside the math: why hexagons}, titlerule=0mm, toprule=.15mm, toptitle=1mm]

The hex bin specifically uses hexagons rather than squares or triangles.
The reason is geometric: among regular polygons that tile the plane, the
hexagon has the most sides, which makes each cell closer to a circle.

For a square grid, a point near the corner of a cell is much further
from the cell center than a point near the middle of an edge. This
produces visible artifacts: the eye sees the grid even when the
underlying data has no gridded structure. Hexagons reduce the effect
dramatically. The visual texture of a hex bin chart is closer to the
underlying density than any rectangular alternative.

This is the same reason bee hives, soap bubbles, and basalt columns are
hexagonal: it is the most efficient regular tiling. Nature picks
hexagons when packing matters; we pick them for the same reason in
visualization.

\end{tcolorbox}

\begin{tcolorbox}[enhanced jigsaw, arc=.35mm, bottomrule=.15mm, bottomtitle=1mm, breakable, colback=white, colbacktitle=quarto-callout-caution-color!10!white, colframe=quarto-callout-caution-color-frame, coltitle=black, left=2mm, leftrule=.75mm, opacityback=0, opacitybacktitle=0.6, rightrule=.15mm, title=\textcolor{quarto-callout-caution-color}{\faFire}\hspace{0.5em}{Reflect before moving on}, titlerule=0mm, toprule=.15mm, toptitle=1mm]

Three questions to articulate in your own words before moving on:

\begin{enumerate}
\def\labelenumi{\arabic{enumi}.}
\item
  Pick one chart type from this chapter you had never made before. What
  question does it answer that the standard repertoire could not?
\item
  Bar charts are the default for almost every quantitative comparison.
  Name a situation where you have used a bar chart that would have been
  better as a Cleveland dot plot.
\item
  The chapter argued that ``the right tool for the question'' is the
  design move. Where do you think your own default chart vocabulary is
  too small?
\end{enumerate}

\end{tcolorbox}

\begin{tcolorbox}[enhanced jigsaw, arc=.35mm, bottomrule=.15mm, bottomtitle=1mm, breakable, colback=white, colbacktitle=quarto-callout-important-color!10!white, colframe=quarto-callout-important-color-frame, coltitle=black, left=2mm, leftrule=.75mm, opacityback=0, opacitybacktitle=0.6, rightrule=.15mm, title=\textcolor{quarto-callout-important-color}{\faExclamation}\hspace{0.5em}{Frontier note}, titlerule=0mm, toprule=.15mm, toptitle=1mm]

The visualization literature has produced many more chart types than the
six in this chapter. Some that are gaining traction in modern empirical
work:

\begin{itemize}
\tightlist
\item
  \emph{Parallel coordinates} for high-dimensional data, showing many
  variables at once with one line per observation
\item
  \emph{Heatmap matrices with hierarchical clustering} for genomics and
  exploratory matrix data
\item
  \emph{Network and graph layouts} for relational data (which the trade
  panel of this book is, when treated as a directed graph)
\item
  \emph{Calendar heatmaps} for time-of-day, day-of-week patterns
\end{itemize}

The next chapter (visualizing uncertainty) covers a separate frontier:
chart types that show not just the estimate but the uncertainty around
it. Most of those are also underused outside of statistical specialty
fields, even though they are exactly what the credibility revolution in
empirical research is asking for.

\end{tcolorbox}

\section{Practice problems}\label{practice-problems-6}

\begin{enumerate}
\def\labelenumi{\arabic{enumi}.}
\item
  Take any spreadsheet of two-time-point data you have access to. Render
  it as a slope graph. Count the number of lines that go up versus down.
  What story does the slope graph tell that you could not have read from
  the means alone?
\item
  The Cleveland dot plot uses position rather than length to encode
  magnitude. Design a dot plot for the data in the trade slice (USA's
  exports to its five partners in 2023). Then redesign it for export
  \emph{change} between 2014 and 2023. What changes between the two
  charts?
\item
  Pick a dataset with more than 10,000 rows and two continuous variables
  (any source: a research dataset, a public CSV, a log file). Make a
  regular scatter plot. Then make a hex-bin version. Compare what each
  chart reveals.
\item
  Find a contingency table in a recent paper or news article. Sketch it
  as a mosaic plot. What relationship between the variables becomes
  visible in the mosaic that was hidden in the table of numbers?
\item
  Construct an alluvial diagram for a categorical transition you
  actually care about (your students' major changes, your team's role
  transitions, partner shifts in a relationship dataset). The exercise
  is partly about whether the data lets you build it cleanly.
\end{enumerate}

\chapter{Visualizing uncertainty}\label{visualizing-uncertainty}

\section{What error bars cannot do}\label{what-error-bars-cannot-do}

Most published research that visualizes uncertainty uses error bars. A
bar chart with a thin vertical line on top of each bar, terminated in
horizontal caps, encoding either the standard error or a 95\% confidence
interval. Error bars are everywhere. They are also, in most situations,
the wrong choice.

Error bars fail at several things at once. They are visually dominated
by the bar; the eye sees the bar's height and treats the error as
decoration. They cannot show asymmetric uncertainty. They cannot show
non-Gaussian uncertainty. They invite the wrong reading: that the
interval is ``where the true value probably is,'' which is closer to a
Bayesian credible interval than to what a frequentist confidence
interval actually means. They fail catastrophically when the bars are
themselves negative, because the visual implication of an ``error bar on
a bar chart'' is not designed to handle bidirectional uncertainty around
a central estimate that swings sign.

This chapter is about better tools. Each one shows uncertainty as a
first-class part of the visual story rather than an afterthought.

\section{Confidence bands and prediction
bands}\label{confidence-bands-and-prediction-bands}

When the underlying quantity is a curve (a regression line, a smoothed
trend, a fitted survival curve), uncertainty is shown as a band around
the curve. The band has a center line (the point estimate) and a shaded
region around it (the uncertainty interval at each point along the
x-axis).

\begin{figure*}

\centering{

\pandocbounded{\includegraphics[keepaspectratio]{prose/part1/../../pictures/figures/output/fig_conf_band.pdf}}

}

\caption{\label{fig-conf-band}A linear regression of bilateral exports
on importer GDP, with a 95\% confidence band shaded around the fitted
line. The band narrows in the dense part of the data and widens at the
extremes, which is a faithful encoding of how much the regression line
itself is uncertain at each x-value. The wider band beyond the data
range is the most honest way to show extrapolation uncertainty.}

\end{figure*}%

Two important distinctions:

A \emph{confidence band} shows the uncertainty around the regression
line itself: where the true mean function plausibly lies. A
\emph{prediction band} shows the uncertainty around future observations:
where new data points are likely to fall. The prediction band is always
wider, because it includes both the uncertainty in the line and the
residual scatter around the line. Confusing them is one of the most
common errors in applied empirical work.

Confidence bands handle asymmetry naturally. If the true function has
heteroskedastic errors, the band can fan out at the edges. They also
handle non-linear models easily: a logistic regression's confidence band
on the probability scale is non-symmetric, and the band reflects this.
They are the right default for any continuous predictor.

\section{Coefficient plots}\label{coefficient-plots}

A coefficient plot, also called a forest plot in medicine and
meta-analysis, shows many regression estimates at once. Each row is one
coefficient. The dot is the point estimate. A horizontal line through
the dot is the confidence interval. A vertical reference line, usually
at zero, marks the null.

\begin{figure}

\centering{

\pandocbounded{\includegraphics[keepaspectratio]{prose/part1/../../pictures/figures/output/fig_coef_plot.pdf}}

}

\caption{\label{fig-coef-plot}Coefficients from a regression of log
bilateral exports on log importer GDP, log distance, common-language
indicator, and an intercept. Each row is one coefficient with a 95\%
interval. The vertical line at zero marks the null. Coefficients whose
intervals cross the line cannot reject the null at the conventional
level; coefficients whose intervals sit clearly to one side are far from
zero in their estimated direction.}

\end{figure}%

Why this works better than the standard alternative, a regression table.
A regression table displays the same information as numbers in a grid.
The numbers require translation: the reader must compare each
coefficient to its standard error mentally, decide whether the interval
crosses zero, and judge magnitudes by mentally normalizing to the units
of each variable. A coefficient plot does all of this visually. The
reader sees at a glance which estimates are large, which are precise,
which cross zero, and which are similar in magnitude.

Coefficient plots are now standard in modern empirical economics and
political science papers, partly because of the credibility revolution
and partly because they make it harder to hide weak results in a table
of numbers.

\section{Hypothetical outcome plots}\label{hypothetical-outcome-plots}

A hypothetical outcome plot, or HOP, shows uncertainty by sampling.
Instead of drawing a confidence band, the plot shows several plausible
realizations of the underlying quantity, drawn from the posterior or the
sampling distribution. In a static figure, an HOP overlays a few dozen
sampled curves in light color, with the central estimate in a heavier
line. In a dynamic figure (animated, interactive), the HOP cycles
through draws over time, giving the reader a visceral sense of how much
the estimate could vary.

\begin{figure*}

\centering{

\pandocbounded{\includegraphics[keepaspectratio]{prose/part1/../../pictures/figures/output/fig_hop.pdf}}

}

\caption{\label{fig-hop}Twenty regression lines drawn from the joint
distribution of slope and intercept, overlaid on the data. The dark line
is the central estimate. Each light line is one plausible realization of
the underlying line consistent with the data. The thicker the visible
bundle, the less confident the estimate; the spread of the bundle is the
uncertainty.}

\end{figure*}%

The pedagogical case for HOPs is that confidence bands feel like static
facts, whereas HOPs feel like the underlying probabilistic process.
Research by Hullman, Kale, and others has shown that HOPs improve
viewers' calibration of subjective beliefs about uncertainty better than
equivalent confidence-band displays. The reason is roughly that bands
encode a precise probability, but the eye reads ``more probable here,
less probable here'' with mixed accuracy, whereas HOPs encode the same
probabilities by the \emph{frequency of visual presence}, which the eye
reads directly.

\section{Fan charts}\label{fan-charts}

A fan chart is a forecast visualization. The forecasted quantity is
plotted as a curve into the future, with a series of nested confidence
bands around it. The innermost band might be a 50\% prediction interval,
the next 80\%, the outermost 95\%. The result looks like a fan opening
to the right, narrowest at the present moment and widening into the
future.

\begin{figure}

\centering{

\pandocbounded{\includegraphics[keepaspectratio]{prose/part1/../../pictures/figures/output/fig_fan_chart.pdf}}

}

\caption{\label{fig-fan-chart}A fan chart of a forecasted series. The
dark central line is the point forecast; the nested bands are the 50\%,
80\%, and 95\% prediction intervals. The fan widens further into the
future, encoding the rate at which forecast uncertainty grows. Bank of
England inflation reports use exactly this format; the visual move is to
communicate both the central forecast and the rate of uncertainty growth
in one figure.}

\end{figure}%

Fan charts are standard in central banking and in epidemiology
forecasting. They communicate two things at once: the central forecast
(the line in the middle) and the rate at which uncertainty grows (the
rate at which the fan opens).

The visual property worth noting: the bands are nested rather than
separate. The reader's eye reads ``more probability concentrated in the
inner bands, less in the outer ones,'' which is the right reading. A
single 95\% band gives a binary in-or-out reading. A nested fan gives a
graded reading: more or less likely.

\section{Showing all the data}\label{showing-all-the-data}

The most honest uncertainty visualization is showing all the data. When
sample sizes permit, plot every individual observation as a small mark,
with the summary (mean, regression line, model) overlaid. The eye reads
both the structure and the spread, and there is no chance of being
misled about precision.

\begin{figure}

\centering{

\pandocbounded{\includegraphics[keepaspectratio]{prose/part1/../../pictures/figures/output/fig_show_all.pdf}}

}

\caption{\label{fig-show-all}Every one of our 300 bilateral observations
plotted as a point, with a smoothed conditional mean overlaid. The
dispersion of the points around the line is the residual variation. The
slope of the line is the estimated relationship. Showing the data and
the model together is the most informationally honest visualization of
uncertainty available.}

\end{figure}%

This is what scatter plots already do, but the principle generalizes. A
box plot shows summary statistics; a strip plot shows summary statistics
and every individual observation. A bar chart shows means; a violin or
jittered scatter shows means and the underlying spread.

The trade-off. When \(n\) is small, ``show all the data'' is easy. When
\(n\) is millions, you need to pre-aggregate or use density-based
encodings (hex bins, contour plots). The principle stays the same:
prefer showing more rather than less, and make explicit the choice about
what is hidden.

\begin{tcolorbox}[enhanced jigsaw, arc=.35mm, bottomrule=.15mm, bottomtitle=1mm, breakable, colback=white, colbacktitle=quarto-callout-warning-color!10!white, colframe=quarto-callout-warning-color-frame, coltitle=black, left=2mm, leftrule=.75mm, opacityback=0, opacitybacktitle=0.6, rightrule=.15mm, title=\textcolor{quarto-callout-warning-color}{\faExclamationTriangle}\hspace{0.5em}{Key insight}, titlerule=0mm, toprule=.15mm, toptitle=1mm]

Uncertainty deserves the same care as the point estimate. Error bars are
common but rarely optimal: they hide more than they show. Confidence
bands, coefficient plots, hypothetical outcome plots, and fan charts
each show uncertainty as a first-class part of the visual argument. The
general rule: show all the data when you can; where summary is
unavoidable, encode the uncertainty in a way that matches the reader's
question.

\end{tcolorbox}

\begin{tcolorbox}[enhanced jigsaw, arc=.35mm, bottomrule=.15mm, bottomtitle=1mm, breakable, colback=white, colbacktitle=quarto-callout-tip-color!10!white, colframe=quarto-callout-tip-color-frame, coltitle=black, left=2mm, leftrule=.75mm, opacityback=0, opacitybacktitle=0.6, rightrule=.15mm, title=\textcolor{quarto-callout-tip-color}{\faLightbulb}\hspace{0.5em}{Try this}, titlerule=0mm, toprule=.15mm, toptitle=1mm]

A regression of bilateral exports on importer GDP gives a slope estimate
of 0.88 with a standard error of 0.06. Sketch:

\begin{enumerate}
\def\labelenumi{\arabic{enumi}.}
\tightlist
\item
  The result as an error-bar plot.
\item
  The result as a coefficient plot, with the null line at 0 and a second
  null line at 1 (the gravity-model benchmark).
\item
  The result as a confidence band over the range of importer GDP in the
  data.
\end{enumerate}

Compare the three. Which most clearly answers ``is the slope different
from 1?'' Which most clearly answers ``how confident am I about the
slope at GDP equal to 25 trillion?'' Which would you publish?

\end{tcolorbox}

\begin{tcolorbox}[enhanced jigsaw, arc=.35mm, bottomrule=.15mm, bottomtitle=1mm, breakable, colback=white, colbacktitle=quarto-callout-note-color!10!white, colframe=quarto-callout-note-color-frame, coltitle=black, left=2mm, leftrule=.75mm, opacityback=0, opacitybacktitle=0.6, rightrule=.15mm, title=\textcolor{quarto-callout-note-color}{\faInfo}\hspace{0.5em}{Inside the math: pointwise vs simultaneous bands}, titlerule=0mm, toprule=.15mm, toptitle=1mm]

A pointwise 95\% confidence band at each \(x\)-value gives an interval
such that, if you repeated the experiment many times, 95\% of the
resulting intervals at that \(x\) would cover the true mean function. A
simultaneous 95\% confidence band is wider and gives an interval such
that the entire band contains the true function 95\% of the time.

These are different things. The pointwise band is the standard default
in most software and is what most researchers report. The simultaneous
band is what you actually want if you intend to make claims about
multiple \(x\)-values at once. The simultaneous band can be
substantially wider than the pointwise version, especially when the
predictor has many evaluation points.

A simple rule of thumb. If you are reporting a single test at a single
\(x\), the pointwise band is fine. If you are reporting on an entire
curve and might select \(x\)-values post hoc, you want the simultaneous
band. The Working-Hotelling band is the classical construction; modern
bootstrap-based bands are the practical alternative.

\end{tcolorbox}

\begin{tcolorbox}[enhanced jigsaw, arc=.35mm, bottomrule=.15mm, bottomtitle=1mm, breakable, colback=white, colbacktitle=quarto-callout-caution-color!10!white, colframe=quarto-callout-caution-color-frame, coltitle=black, left=2mm, leftrule=.75mm, opacityback=0, opacitybacktitle=0.6, rightrule=.15mm, title=\textcolor{quarto-callout-caution-color}{\faFire}\hspace{0.5em}{Reflect before moving on}, titlerule=0mm, toprule=.15mm, toptitle=1mm]

Three questions to articulate in your own words:

\begin{enumerate}
\def\labelenumi{\arabic{enumi}.}
\item
  The chapter claims that error bars on bar charts are usually the wrong
  choice. In your own words, why? Where have you seen error bars used
  well?
\item
  A confidence band and a prediction band are different. Which one would
  you want to see if you were asked to forecast where a new observation
  will land? Which if you were asked about the underlying mean function?
\item
  HOPs (hypothetical outcome plots) work because the eye reads frequency
  directly. Where else in this book have you seen frequency used as a
  visual encoding?
\end{enumerate}

\end{tcolorbox}

\begin{tcolorbox}[enhanced jigsaw, arc=.35mm, bottomrule=.15mm, bottomtitle=1mm, breakable, colback=white, colbacktitle=quarto-callout-important-color!10!white, colframe=quarto-callout-important-color-frame, coltitle=black, left=2mm, leftrule=.75mm, opacityback=0, opacitybacktitle=0.6, rightrule=.15mm, title=\textcolor{quarto-callout-important-color}{\faExclamation}\hspace{0.5em}{Frontier note}, titlerule=0mm, toprule=.15mm, toptitle=1mm]

Modern uncertainty visualization is an active research area. Three
directions worth knowing.

\emph{Animated and interactive uncertainty.} Hypothetical outcome plots,
ensemble visualizations, and dynamic confidence regions have been shown
by Hullman, Kale, and others to improve viewer calibration of subjective
uncertainty. The visualizations on this book's website use these where
they fit.

\emph{Multiverse analysis displays.} When the same dataset can support
many defensible analyses (specification choices, measurement choices,
robustness checks), the modern move is to display the distribution of
results across all of them. Specification curves and multiverse plots
are the two leading visual formats.

\emph{Bayesian posterior visualization.} Modern Bayesian software
produces full posterior distributions. Visualizing those distributions,
rather than reducing them to a single CI, is its own active design
space. Density plots, eye-plots, and fan-style visualizations are all in
current use.

The general direction: less reduction, more information. The trend
follows the technical capacity to compute and display the underlying
probabilistic objects directly, rather than summarize them away.

\end{tcolorbox}

\section{Practice problems}\label{practice-problems-7}

\begin{enumerate}
\def\labelenumi{\arabic{enumi}.}
\item
  Find a paper in your field that uses error bars on a bar chart. Sketch
  the same data as a coefficient plot. Which version is more useful for
  the question the paper is trying to answer?
\item
  A 95\% confidence band for a regression of \(y\) on \(x\) is plotted.
  At \(x = 5\) the band runs from \(y = 2.4\) to \(y = 3.6\). Write one
  sentence describing what this band means, and one sentence describing
  what it does \emph{not} mean.
\item
  Take the trade slice. Fit a simple regression of log exports on log
  importer GDP. Plot it three ways: (a) point estimate plus error bar;
  (b) confidence band; (c) hypothetical outcome plot with twenty draws.
  Compare. Which makes the gravity slope feel most uncertain?
\item
  A fan chart of disease cases shows the 50\%, 80\%, and 95\% prediction
  intervals widening into the future. Two months out, the 95\% interval
  spans 100 to 1,000 cases. Is that a useful forecast? When would a fan
  chart with very wide outer bands still be informative?
\item
  Construct a coefficient plot for the multiple-regression result you
  choose: at minimum three coefficients with their 95\% intervals and a
  vertical reference line at zero. Annotate which coefficients are
  statistically distinguishable from zero and which are not. The
  exercise is half the chart and half the writing of the annotations.
\end{enumerate}


\backmatter


\end{document}
